% !TEX program = /Library/TeX/texbin/latexmk
% !TEX options = -pdf -synctex=1 -interaction=nonstopmode -halt-on-error -file-line-error "%DOC%"
\documentclass[11pt,a4paper]{article}
\usepackage[T1]{fontenc}
\usepackage[utf8]{inputenc}
\usepackage{lmodern}
\usepackage[a4paper,margin=21mm,headheight=15pt]{geometry}
\usepackage{microtype}
\usepackage{amsmath,amssymb}
\usepackage{booktabs,longtable,array,tabularx}
\usepackage{needspace}
\usepackage{xcolor}
\usepackage{listings}
\usepackage{enumitem}
\usepackage{tikz}
\usetikzlibrary{arrows.meta,positioning}
\usepackage{fancyhdr}
\usepackage{hyperref}
\usepackage{xurl}
\definecolor{navy}{HTML}{17324D}
\definecolor{teal}{HTML}{176B72}
\definecolor{pale}{HTML}{F3F6F8}
\definecolor{muted}{HTML}{546473}
\hypersetup{colorlinks=true,linkcolor=navy,citecolor=teal,urlcolor=teal,
 pdftitle={Detailed code reviews of DWD simulation, SNR and subtraction},
 pdfauthor={Codex},bookmarksnumbered=true}
\pagestyle{fancy}
\fancyhf{}
\fancyhead[L]{\small\sffamily DWD workflows: detailed code review}
\fancyhead[R]{\small\sffamily 13 October 2026}
\fancyfoot[C]{\small\thepage}
\renewcommand{\headrulewidth}{0.3pt}
\setlength{\parindent}{0pt}
\setlength{\parskip}{5pt}
\setlength{\emergencystretch}{3em}
\clubpenalty=10000
\widowpenalty=10000
\setcounter{tocdepth}{2}
\makeatletter
\renewcommand{\l@subsection}{\@dottedtocline{2}{1.5em}{3.2em}}
\makeatother
\setlist[itemize]{leftmargin=1.6em,itemsep=2pt,topsep=4pt}
\setlist[enumerate]{leftmargin=1.8em,itemsep=3pt,topsep=4pt}
\newcommand{\code}[1]{\texttt{\nolinkurl{#1}}}
\newcommand{\linesref}[2]{\textbf{Lines #1--#2}}
\newcommand{\mainfile}{\code{eric_subtraction_all.py}}
\newcommand{\Hz}{\,\mathrm{Hz}}
\newcommand{\dd}{\,\mathrm{d}}
\newcommand{\ip}[2]{\left(#1\,\middle|\,#2\right)}
\newcommand{\status}[1]{\textbf{\textcolor{teal}{#1}}}
\newcolumntype{P}[1]{>{\raggedright\arraybackslash}p{#1}}
\lstdefinestyle{reviewpython}{
 language=Python,basicstyle=\ttfamily\footnotesize,
 keywordstyle=\color{navy}\bfseries,commentstyle=\color{muted},
 stringstyle=\color{teal},numbers=left,stepnumber=1,numberblanklines=true,showlines=false,numberstyle=\tiny\color{muted},
 numbersep=7pt,frame=single,rulecolor=\color{black!20},
 backgroundcolor=\color{pale},breaklines=true,breakatwhitespace=false,
 showstringspaces=false,columns=fullflexible,keepspaces=true,tabsize=4,
 xleftmargin=13pt,aboveskip=8pt,belowskip=8pt}
\lstset{style=reviewpython}



\begin{document}
\hypersetup{pageanchor=false}
\begin{titlepage}
\vspace*{6mm}
\begin{center}
{\sffamily\bfseries\color{navy}\Huge \resizebox{\textwidth}{!}{DWD injection and subtraction workflow}\par}
\vspace{5mm}
{\sffamily\Large Detailed code reviews of simulation, SNR and subtraction\par}
\vspace{4mm}
\end{center}
\vspace{4mm}


\begin{abstract}
Hallowed is the labour that hearkens to the hymn of double white dwarfs. Though their ancient fires have faded, their quiet witness endures: two pale embers bearing their song across the darkness. Blessed are the hands that seek their hidden harmony with patience and humility. For in their celestial dance we glimpse the deep thoughts of the Creator, borne upon the waves of the unseen.
\end{abstract}


\end{titlepage}
\hypersetup{pageanchor=true}
\pagenumbering{roman}
\begingroup
\small\setlength{\parskip}{0pt}
\tableofcontents
\endgroup
\clearpage
\pagenumbering{arabic}

\Needspace{8\baselineskip}
\section{DWD simulation pipeline}
\label{sec:simulation_pipeline}

The subtraction workflow configured in Section~\ref{sec:walkthrough1} operates on a precomputed DWD observation. This section describes how \code{eric_simulation_DWD.py} produces that input through population injection and detector simulation.

\Needspace{10\baselineskip}
\Needspace{8\baselineskip}
\subsection{Lines 1--49: import packages and configure the simulation}
% SOURCE eric_simulation_DWD.py 1 20
\begin{lstlisting}[language=Python,firstnumber=1,numbers=left]
# Step 1. Import packages.
import os
for thread_variable in ('OMP_NUM_THREADS', 'OPENBLAS_NUM_THREADS', 'MKL_NUM_THREADS', 'NUMEXPR_NUM_THREADS'):
    os.environ[thread_variable] = '1'
import importlib.metadata
import json
from pathlib import Path
import shutil
import tempfile
import time
import h5py
import lisaorbits
import lisagwresponse
import lisainstrument
import numpy as np
import pytdi
from script.oneblock_observation import pack_observation
from eric_generate_waveform import DWD_simulation_gpu, tdi_gpu


\end{lstlisting}

Import all the necessary packages.

% SOURCE eric_simulation_DWD.py 21 37
\begin{lstlisting}[language=Python,firstnumber=21,numbers=left]
# Step 2. Configure parameter settings.
catalogue_file = Path('./data/LDC2_sangria_training_v2.h5')
observation_file = Path('./data/dwd_observation_full_catalogue.h5')
T_obs = 360. * 86400.
fs = 0.1
L = 2.5e9
N_acc = 3.0e-15
delta_x = 15.0e-12
gpu_batch_size = 10000
minimum_envelope_samples = 512
parameter_fields = dict(
    A='Amplitude', f='Frequency', df='FrequencyDerivative', phi0='InitialPhase',
    iota='Inclination', longitude='EclipticLongitude', latitude='EclipticLatitude',
    psi_bme='Polarization',
)


\end{lstlisting}

The two paths specify the source catalogue and combined observation file. \code{T_obs}, \code{fs}, \code{L}, \code{N_acc} and \code{delta_x} set the observation duration, sampling frequency, arm length, acceleration-noise level and OMS-noise level.

\code{gpu_batch_size=10000} limits the number of ordinary DWDs processed together on the GPU. \code{minimum_envelope_samples=512} sets the minimum waveform-envelope resolution; wider chirps use smaller batches. The dictionary maps the eight internal parameter names to their catalogue fields.

% SOURCE eric_simulation_DWD.py 38 49
\begin{lstlisting}[language=Python,firstnumber=38,numbers=left]
workspace = Path(tempfile.mkdtemp(prefix='.dwd_simulation_', dir=observation_file.parent))
orbit_file = workspace / 'orbit.h5'
N = int(round(T_obs * fs))
N_day = int(round(86400 * fs))
detector_config = dict(
    T_obs=T_obs, gpu_batch_size=gpu_batch_size,
    minimum_envelope_samples=minimum_envelope_samples, include_padding=True,
    dt=1 / fs, orbits=str(orbit_file), orbit_dataset='tcb/ltt',
    aafilter=('kaiser', 240., .275 * fs, .725 * fs),
    clock_offsets=0., clock_freqoffsets=0., clock_freqlindrifts=0., clock_freqquaddrifts=0.,
)

\end{lstlisting}

The temporary workspace holds intermediate files until publication. \code{N} is the observation sample count, and \code{N_day} supplies one day of boundary padding. The single \code{detector_config} dictionary supplies both GPU injection and the native instrument.

\Needspace{8\baselineskip}
\subsection{Lines 50--66: create the orbit and continuous noise}
% SOURCE eric_simulation_DWD.py 50 66
\begin{lstlisting}[language=Python,firstnumber=50,numbers=left]
# Step 3. Simulate instrumental noise with the LISA packages.
lisaorbits.KeplerianOrbits(L=L).write(
    str(orbit_file), t0=-20e6, dt=86400., size=825, mode='w',
)
print('Simulating one continuous ACC+OMS noise realization.', flush=True)
noise_instrument = lisainstrument.Instrument(
    **{key: value for key, value in detector_config.items()
       if key not in ('T_obs', 'gpu_batch_size', 'minimum_envelope_samples', 'include_padding')}, size=N + 2 * N_day,
    t0=-N_day / fs,
    testmass_asds=N_acc, oms_asds=(delta_x, 0., 0., 0., 0., 0.),
)
noise_instrument.disable_all_noises(excluding={'test-mass', 'oms'})
noise_instrument.disable_dopplers()
noise_data = pytdi.Data.from_instrument(
    noise_instrument.export_numpy_full(),
)

\end{lstlisting}

The orbit table starts at \code{t0=-20e6}, uses one-day sampling and contains 825 samples, providing coverage for interpolation and delayed measurements. \code{noise_instrument} produces one continuous ACC+OMS measurement realization with \(N+2N_{\rm day}\) samples beginning at \(-N_{\rm day}/f_s\), retaining one padding day at each end. \code{pytdi.Data.from_instrument} exposes the time-domain noise measurements for later combination.

\Needspace{10\baselineskip}
\Needspace{8\baselineskip}
\subsection{Lines 67--112: prepare the catalogue and full-duration signal measurements}
% SOURCE eric_simulation_DWD.py 67 83
\begin{lstlisting}[language=Python,firstnumber=67,numbers=left]
# Step 4. Simulate waveforms of all DWDs on the GPU.
with h5py.File(catalogue_file) as source:
    dgb = source['sky/dgb/cat'][:]
    igb = source['sky/igb/cat'][:]

source_count = len(dgb) + len(igb)
catalogue = np.empty(source_count, dtype=[('ID', 'i8')] + [
    (field, 'f8') for field in parameter_fields.values()
])
catalogue['ID'] = np.arange(1, source_count + 1, dtype=np.int64)
for field in parameter_fields.values():
    catalogue[field] = np.concatenate((dgb[field][:, 0], igb[field][:, 0]))
del dgb, igb

with h5py.File(workspace / 'population.h5', 'x') as saved:
    saved.create_dataset('catalogues/all', data=catalogue, chunks=True)

\end{lstlisting}

All catalogue rows are combined into one table with a unique ID and eight physical parameters per DWD. This table is saved for the subsequent injection and subtraction analyses.

% SOURCE eric_simulation_DWD.py 84 99
\begin{lstlisting}[language=Python,firstnumber=84,numbers=left]
links_file = workspace / 'population_links.h5'
print(f'Simulating all {source_count:,} DWDs with DWD_simulation_gpu.', flush=True)
started = time.perf_counter()
links = DWD_simulation_gpu(catalogue, detector_config)
waveform_runtime = time.perf_counter() - started
print(f'GPU injection completed: {source_count:,} DWDs in {waveform_runtime:.1f}s.', flush=True)
del catalogue
response_model = lisagwresponse.GalacticBinary(
    A=1., f=.003, ra=0., dec=0., orbits=str(orbit_file),
    orbit_interp_order=5, shift_sun='never',
)
with h5py.File(links_file, 'x') as saved:
    response_model.init_response_file(saved, detector_config['dt'], len(links), -2 * N_day / fs, 'tcb')
    saved['tcb/y'][:] = links
del links

\end{lstlisting}

\mbox{\code{DWD_simulation_gpu}} receives the DWD parameter table and detector settings, then returns the summed six-link time-domain response. It accepts the eight physical parameters directly, so sources need not belong to the LDC catalogue; any ID column is retained only as table metadata. Lines 91--97 store the summed GPU response in the format required by \code{lisainstrument}.

% SOURCE eric_simulation_DWD.py 100 112
\begin{lstlisting}[language=Python,firstnumber=100,numbers=left]
# Convert the summed links into full-duration detector measurements.
signal_instrument = lisainstrument.Instrument(
    **{key: value for key, value in detector_config.items()
       if key not in ('T_obs', 'gpu_batch_size', 'minimum_envelope_samples', 'include_padding')}, gws=str(links_file), size=N + 2 * N_day,
    t0=-N_day / fs,
)
signal_instrument.disable_all_noises()
signal_instrument.disable_dopplers()
signal_data = pytdi.Data.from_instrument(signal_instrument.export_numpy_full())
signal_measurements = signal_data.measurements
noise_measurements = noise_data.measurements
del noise_data, noise_instrument

\end{lstlisting}

The summed links are converted into detector measurements over the entire padded observation in one call. At the end of this subsection, \code{signal_measurements} and \code{noise_measurements} are ready on the same time grid.

\Needspace{8\baselineskip}
\subsection{Lines 113--134: combine the measurements, apply TDI, and save}
\label{sec:signal_noise_tdi}
% SOURCE eric_simulation_DWD.py 113 134
\begin{lstlisting}[language=Python,firstnumber=113,numbers=left]
# Step 5. Combine the full observation, apply TDI, and save.
total_measurements = {
    name: signal_measurements[name] + noise_measurements[name]
    for name in signal_measurements
}
tdi_results = {}
for name, measurements in (('total', total_measurements),
                           ('signal', signal_measurements), ('noise', noise_measurements)):
    tdi_results[name] = tdi_gpu(
        measurements, signal_data.args, keep=slice(N_day, N_day + N),
    ) / signal_instrument.central_freq
detector_config['orbits'] = str(observation_file.with_suffix('.orbits.h5'))
with h5py.File(workspace / 'population_tdi.h5', 'x') as saved:
    saved.attrs.update(duration=T_obs, rate=fs)
    saved.attrs['detector_config'] = json.dumps(detector_config)
    for name, value in tdi_results.items():
        saved.create_dataset(name + '_xyz_td', data=value, chunks=(3, min(N_day, N)))

pack_observation(workspace, observation_file, overwrite=True, time_domain_only=True)
shutil.copy2(workspace / 'orbit.h5', observation_file.with_suffix('.orbits.h5'))
shutil.rmtree(workspace)
print(f'Saved {source_count:,} DWDs: {observation_file.resolve()}', flush=True)
\end{lstlisting}

The matching signal and noise measurements are added before TDI. \code{tdi_gpu} applies the installed PyTDI \code{X2/Y2/Z2} expressions with GPU interpolation, reusing the moving-arm delays for the total, signal and noise streams. It retains the native fifth-order delay interpolation and 31st-order measurement interpolation; the implementation and comparison with standard PyTDI are given in Section~\ref{sec:tdi_gpu_api} and Section~\ref{sec:tdi_api_validation}. The file stores only the final time-domain TDI arrays for total, signal and noise, together with the DWD table, orbit and detector settings. The observation also stores \code{duration=T_obs}, \code{rate=fs}, and the complete \code{detector_config} as JSON. The subtraction pipeline reads these saved values directly instead of redefining detector settings. The native orbit is also exported as a companion \code{.orbits.h5} file for the subtraction pipeline, preserving the orbit used to simulate the observation.


\clearpage
\Needspace{8\baselineskip}
\section{Detailed code review of eric\_snr.py}
\label{sec:snr_pipeline}

This section reproduces all 360 lines of \code{eric_snr.py} in source order.
Its four numbered stages prepare the SNR catalogue used by the subtraction
pipeline. Line numbers here belong to \code{eric_snr.py}; the subsequent
subtraction walkthrough resumes its own numbering.

\Needspace{18\baselineskip}
\subsection{Imports, numerical settings and temporary output (lines 1--48)}
\Needspace{6\baselineskip}
% SOURCE eric_snr.py 1 27
\begin{lstlisting}[language=Python,numbers=left,showlines=true,firstnumber=1]
import os
for variable in ('OMP_NUM_THREADS', 'OPENBLAS_NUM_THREADS', 'MKL_NUM_THREADS', 'NUMEXPR_NUM_THREADS'):
    os.environ[variable] = '1'
import logging
import time
from functools import partial
from pathlib import Path
from tempfile import TemporaryDirectory

import h5py
import jax
import jax.numpy as jnp
import numpy as np
from jax.scipy.linalg import lu_factor, lu_solve
from jax.scipy.signal import fftconvolve
from tqdm.auto import tqdm

from script.eric_generate_waveform import PARAMETER_FIELDS, generate_waveform_td
from script.spectrum import DetectorFunction

jax.config.update('jax_enable_x64', True)

logging.getLogger('lisainstrument.instru.instru_fplan').addFilter(
    lambda record: not record.getMessage().startswith(
        'Using default set of locking beatnote frequencies'))


\end{lstlisting}

\linesref{1}{3} set the listed numerical-library thread counts to one before
NumPy and the other scientific packages are imported. This is an existing
setting in the SNR producer. It controls host-library threads; waveform
generation, FFTs and SNR contractions are dispatched through the GPU path.

\linesref{4}{19} import file handling, timing, progress display, HDF5,
NumPy and JAX. \code{lu_factor} and \code{lu_solve} factor and solve the
three-channel covariance systems. \code{fftconvolve} smooths the foreground
spectrum. The imported \code{generate_waveform_td} is the generator in
\code{script/eric_generate_waveform.py}. It yields individual DWD XYZ batches;
the \code{DWD_simulation_gpu} function used by the observation pipeline
returns summed six-link responses. These are separate interfaces.

Line 21 enables JAX's 64-bit numerical mode. Lines 23--25 suppress one
repeated instrument log message about default locking beatnote frequencies.
The filter changes which messages are displayed; it does not alter the
instrument configuration or the SNR calculation.

\Needspace{6\baselineskip}
% SOURCE eric_snr.py 28 48
\begin{lstlisting}[language=Python,numbers=left,showlines=true,firstnumber=28]
LDC_file = Path('./data/LDC2_sangria_training_v2.h5')
orbit_file = Path('./data/lisa_snr_orbits.h5')
output_file = Path('./data/ldc_dwd_data.h5')
catalogue_groups = ('dgb', 'igb')
T_obs = 360 * 86400
fs = 0.1
N_acc = 3.0e-15
delta_x = 15.0e-12
SNR_threshold = 7.0
covariance_regularization = 1.0e-12  # Add this times trace(S_instrument)/3 to all channels.
catalogue_chunk_size = 65536

started = time.perf_counter()
output_file.parent.mkdir(parents=True, exist_ok=True)
temporary_directory = TemporaryDirectory(prefix='ldc_snr_', dir=output_file.parent)
work_directory = Path(temporary_directory.name)
final_output_file = output_file
output_file = work_directory/final_output_file.name
background_file = work_directory/'background.h5'
history = []

\end{lstlisting}

\linesref{28}{38} define this producer's input paths and scientific settings.
\code{LDC_file} supplies the original DWD parameters, \code{orbit_file}
supplies the orbit to the waveform and covariance helpers, and
\code{output_file} names the completed SNR result. This script does not read
the simulated observation from Section~\ref{sec:simulation_pipeline}.

The observation duration is \(T=360\times86400=31{,}104{,}000\) seconds
and the sampling rate is \(f_s=0.1\Hz\). The two noise amplitudes are
\(3.0\times10^{-15}\,\mathrm{m\,s^{-2}/\sqrt{Hz}}\) for acceleration and
\(15.0\times10^{-12}\,\mathrm{m/\sqrt{Hz}}\) for optical metrology.
\code{SNR_threshold=7.0} controls membership of the unresolved foreground.
The regularization coefficient is applied to the instrumental covariance
in lines 151--152. \code{catalogue_chunk_size=65536} limits the number of
catalogue rows handled by each outer file-processing chunk; the imported
waveform generator divides that chunk into its own GPU batches.

\code{catalogue_groups} names the two storage locations in the original LDC
file. Both are processed by the same waveform, covariance and SNR operations.
Their rows are later written into one numerical \code{catalogue} table.

\linesref{40}{48} begin the runtime measurement and prepare a temporary
directory inside the output's parent. The producer creates the parent if
needed. The completed output path is saved as \code{final_output_file};
the working \code{output_file} and \code{background_file} are then placed
under the temporary directory. \code{history} records one diagnostic row
per foreground round. Publication and cleanup occur at the end.

\Needspace{18\baselineskip}
\subsection{Step 1: read the catalogue and prepare the frequency grid (lines 50--69)}
\Needspace{6\baselineskip}
% SOURCE eric_snr.py 49 68
\begin{lstlisting}[language=Python,numbers=left,showlines=true,firstnumber=49]
# -----------------------------------------------------------------------
# 1. Read every DWD catalogue and prepare the native physical frequency grid.
# -----------------------------------------------------------------------
print('1/4  Read the complete LDC DWD catalogue', flush=True)
# Positive FFT bins, excluding DC and Nyquist.
first_bin = 1
last_bin = int(T_obs*fs)//2-1
frequency = np.arange(first_bin, last_bin+1)/T_obs
catalogue_sizes = {}
with h5py.File(LDC_file, 'r') as ldc:
    for group in catalogue_groups:
        catalogue = ldc[f'sky/{group}/cat']
        catalogue_sizes[group] = len(catalogue)
catalogue_offsets = {'dgb': 0, 'igb': catalogue_sizes['dgb']}
source_count = sum(catalogue_sizes.values())
print(f'     {source_count:,} binaries; one {T_obs/86400:g}-day observation; '
      f'{len(frequency):,} native bins; '
      f'GPU TD waveform batches; GPU FFT/SNR and foreground', flush=True)


\end{lstlisting}

\linesref{54}{56} construct the positive Fourier frequencies
\[
N=\mathrm{int}(Tf_s)=3{,}110{,}400,\qquad
f_k=\frac{k}{T},\qquad k=1,\ldots,\frac{N}{2}-1.
\]
The resulting array contains \(K=1{,}555{,}199\) bins with spacing
\(\Delta f=1/T\simeq3.2150\times10^{-8}\Hz\). DC and the Nyquist bin
are excluded. This grid is determined by the full observation duration;
it is not the subtraction pipeline's indexed frequency-block grid.

\linesref{57}{63} read the lengths of the two original catalogue datasets.
The first dataset starts at global index zero; the second starts after
all rows in the first. These offsets let the waveform generator identify
sources consistently while reading the original file. \code{source_count}
is their sum, currently 29,000,000. The final numerical table uses these
same global row indices plus one as its saved IDs.

\linesref{64}{66} report the source count, duration and frequency count.
They do not start a second observation or divide the time series into days.
Each generated source waveform spans the full \(T\)-second observation.

\Needspace{18\baselineskip}
\subsection{GPU SNR, waveform batches and foreground covariance (lines 69--135)}
\Needspace{6\baselineskip}
% SOURCE eric_snr.py 69 76
\begin{lstlisting}[language=Python,numbers=left,showlines=true,firstnumber=69]
@jax.jit
def SNR_gpu(waveform_fd, covariance_factor, T_obs, fs):
    vectors = waveform_fd.transpose(2, 1, 0)
    weighted = lu_solve(covariance_factor, vectors)
    power = jnp.sum((vectors.conj()*weighted).real, axis=(0, 1))
    return jnp.sqrt(4/(T_obs*fs**2)*power)


\end{lstlisting}

\linesref{69}{74} define the GPU SNR contraction. For a batch containing
\(B\) sources, \code{waveform_fd} has shape \((B,3,K)\). The three entries
are the second-generation XYZ channels. Transposing it gives
\code{vectors} with shape \((K,3,B)\), so all sources can be solved against
the same \(3\times3\) covariance matrix at each frequency.

\code{lu_solve} calculates \(\mathbf S_k^{-1}\mathbf H_b[k]\) using an
already computed LU factorization. The sum in line 73 contracts the channel
and frequency axes, leaving one scalar per source. The real part removes
numerical imaginary residue from the Hermitian quadratic form.

The implemented normalization is
\begin{equation}
\rho_b^2=\frac{4}{Tf_s^2}
\sum_{k=1}^{N/2-1}
\operatorname{Re}\!\left[
\mathbf H_b[k]^\dagger\mathbf S(f_k)^{-1}\mathbf H_b[k]\right].
\label{eq:snr_producer_normalization}
\end{equation}
Here \(H_{bi}[k]=\sum_{n=0}^{N-1}h_{bi}[n]\exp(-2\pi\mathrm{i}kn/N)\)
is the unscaled discrete FFT. Its physical Fourier amplitude is
\(H_{bi}[k]/f_s\), and the frequency spacing is \(1/T\), which gives
the factor \(4/(Tf_s^2)\) for a one-sided noise covariance. Line 74 returns
\(\rho_b\), rather than \(\rho_b^2\). The full covariance includes
cross-channel correlations; XYZ SNRs are not combined as three independent
noise channels.

\Needspace{6\baselineskip}
% SOURCE eric_snr.py 77 93
\begin{lstlisting}[language=Python,numbers=left,showlines=true,firstnumber=77]
class GPUJobs:
    def __init__(self, T_obs, fs, orbit_file, first_bin):
        self.T_obs = T_obs
        self.fs = fs
        self.orbit_file = orbit_file
        self.first_bin = first_bin
        self.noise_input = None
        self.noise_covariance = None
        self.covariance_factor = None
        self.foreground_lines = None

    def snr(self, dwd_ids, noise_covariance, chunk=None):
        return self._results(dwd_ids, noise_covariance, None, chunk)

    def foreground(self, dwd_ids, noise_covariance, chunk=None):
        return self._results(dwd_ids, noise_covariance, SNR_threshold, chunk)

\end{lstlisting}

\linesref{77}{86} define the state shared by waveform batches.
\code{T_obs}, \code{fs}, \code{orbit_file} and \code{first_bin} preserve
the observation setup. The other attributes initially contain
\code{None}: the most recently supplied noise array, its GPU copy,
its LU factorization and the accumulated foreground spectral matrices.

\linesref{88}{92} expose two routes through the same calculation.
\code{snr(...)} passes \code{threshold=None}, so it evaluates SNRs without
adding any source to the foreground accumulator.
\code{foreground(...)} passes \code{SNR_threshold}, so it evaluates the
same SNRs and also adds sources at or below that threshold. Both return
an iterator over completed source batches.

\Needspace{6\baselineskip}
% SOURCE eric_snr.py 94 106
\begin{lstlisting}[language=Python,numbers=left,showlines=true,firstnumber=94]
    @partial(jax.jit, static_argnums=(0, 3), donate_argnums=(4,))
    def _process(self, xyz, covariance_factor, threshold, foreground_lines):
        waveform_fd = jnp.fft.rfft(
            xyz, n=int(round(self.T_obs*self.fs)), axis=-1)
        waveform_fd = waveform_fd[
            :, :, self.first_bin:self.first_bin+len(covariance_factor[0])]
        snr = SNR_gpu(waveform_fd, covariance_factor, self.T_obs, self.fs)
        if threshold is not None:
            unresolved = (snr <= threshold).astype(waveform_fd.real.dtype)
            foreground_lines += 2/(self.T_obs*self.fs**2)*jnp.einsum(
                'bik,bjk,b->kij', waveform_fd, waveform_fd.conj(), unresolved)
        return snr, foreground_lines

\end{lstlisting}

\linesref{94}{105} perform the FFT and optional foreground accumulation.
\code{xyz} has shape \((B,3,N)\). The real FFT acts along time, then
lines 98--99 retain the same \(K\) positive bins prepared in Step 1.
The rectangular full-duration time window is implicit: the code applies
no taper before this FFT.

When \code{threshold} is present, line 102 assigns weight one to
sources with \(\rho_b\leq7\) and zero to the others. The threshold
equality therefore belongs to the unresolved set. The accumulated
one-sided cross-spectral matrix is
\begin{equation}
L_{kij}\leftarrow L_{kij}
+\frac{2}{Tf_s^2}\sum_{b:\,\rho_b\leq7}
H_{bi}[k]H_{bj}[k]^*.
\label{eq:snr_foreground_lines}
\end{equation}
The \code{einsum} expression sums over sources while retaining frequency
and both channel axes, producing shape \((K,3,3)\).
It preserves each source's XYZ correlations. Separate sources contribute
their own spectral matrices; the code does not form inter-source coherent
cross terms by summing all waveforms before taking an FFT.

The JIT decorator treats \code{self} and \code{threshold} as static
arguments. Donating the previous foreground buffer lets JAX reuse its
storage when returning the updated accumulator. The \code{threshold=None}
route skips the update and returns the existing accumulator alongside
the newly calculated SNRs.

\Needspace{6\baselineskip}
% SOURCE eric_snr.py 107 125
\begin{lstlisting}[language=Python,numbers=left,showlines=true,firstnumber=107]
    def _results(self, dwd_ids, noise_covariance, threshold, chunk=None):
        description = 'Generate full TD XYZ'
        if chunk is not None:
            description += f' [{chunk[0]}/{chunk[1]}]'
        for batch_ids, _, xyz in generate_waveform_td(
                self.fs, self.T_obs, dwd_ids, self.orbit_file,
                progress_desc=description):
            if noise_covariance is not self.noise_input:
                self.noise_covariance = jax.device_put(
                    noise_covariance, jax.devices('gpu')[0])
                self.covariance_factor = lu_factor(self.noise_covariance)
                self.noise_input = noise_covariance
            if self.foreground_lines is None:
                self.foreground_lines = jnp.zeros_like(self.noise_covariance)
            snr, self.foreground_lines = self._process(
                xyz, self.covariance_factor, threshold, self.foreground_lines)
            jax.block_until_ready((snr, self.foreground_lines))
            yield batch_ids, np.asarray(snr)

\end{lstlisting}

\linesref{107}{124} connect the source IDs to the waveform and SNR
calculations. The imported generator yields \code{batch_ids}, sample
times and \code{xyz}; this caller discards the sample-time array because
\(T\), \(f_s\) and the FFT grid are already fixed. In the inspected
generator, the fast route processes batches of up to 64 sources, with
shape \((B,3,N)\). That internal batch size is distinct from the
65,536-row catalogue chunk.

Lines 114--118 reuse the GPU noise array and its LU factors while the
caller supplies the same covariance object. A different covariance
object, such as the next round's effective covariance, is transferred
and factorized again. The code uses object identity for this cache;
it does not compare array values.

Lines 119--120 initialize the foreground accumulator once with shape
\((K,3,3)\). It remains accumulated across catalogue chunks and subsequent
foreground-transfer rounds. Line 123 waits for both returned GPU arrays
to finish before advancing progress. Line 124 copies only the batch SNR
values to a NumPy array for catalogue bookkeeping; the covariance and
foreground arrays remain on the GPU.

This interface passes source IDs to the imported generator. The inspected
generator reads their parameters from the project's
\code{data/LDC2_sangria_training_v2.h5}; the current default
\code{LDC_file} names that same catalogue. The generator does not receive
\code{LDC_file} as an argument, so its catalogue path must remain aligned
with the file used by this producer's row-reading code.

The generator is passed the specified orbit file for the time-dependent
source response. The SNR helper does not replace that moving response
with a single-frequency template.

\Needspace{6\baselineskip}
% SOURCE eric_snr.py 126 135
\begin{lstlisting}[language=Python,numbers=left,showlines=true,firstnumber=126]
    def confusion_covariance(self):
        with jax.default_device(self.foreground_lines.device):
            offsets = jnp.arange(-128, 129)
            kernel = 0.5*(1+jnp.cos(jnp.pi*offsets/128))
            kernel /= kernel.sum()
            covariance = fftconvolve(
                self.foreground_lines, jnp.broadcast_to(kernel[:, None, None],
                    (len(kernel), 3, 3)), mode='same', axes=0)
        return 0.5*(covariance+covariance.conj().swapaxes(-1, -2))

\end{lstlisting}

\linesref{126}{134} convert the accumulated spectral lines into a
smoothed unresolved-foreground covariance. The normalized kernel is
\[
w_m=\frac{\tfrac12[1+\cos(\pi m/128)]}
{\sum_{q=-128}^{128}\tfrac12[1+\cos(\pi q/128)]},
\qquad m=-128,\ldots,128.
\]
There are 257 kernel positions, with zero weight at the endpoints.
The nominal smoothing half-width is \(128/T\simeq4.1152\times10^{-6}\Hz\).

The convolution runs along frequency only and applies the same kernel to
each of the nine XYZ covariance entries. \code{mode='same'} keeps the
output shape \((K,3,3)\); its finite-array edges use the convolution's
zero extension without a separate edge-weight renormalization.
Line 134 explicitly forms
\(\mathbf C=(\mathbf C+\mathbf C^\dagger)/2\), preserving a Hermitian
matrix at every frequency.

The result is a spectral covariance model for unresolved DWDs.
It is not a noise time series and is not a residual produced by fitting
and subtracting sources from the simulated observation.

\Needspace{18\baselineskip}
\subsection{Step 2: instrumental covariance and absolute SNRs (lines 137--255)}
\Needspace{6\baselineskip}
% SOURCE eric_snr.py 136 154
\begin{lstlisting}[language=Python,numbers=left,showlines=true,firstnumber=136]
# -----------------------------------------------------------------------
# 2. Instrument-only XYZ covariance, then absolute SNRs for every source.
# -----------------------------------------------------------------------
print('2/4  Instrumental covariance and round 1 absolute SNRs', flush=True)
t = T_obs/2
with tqdm(total=1, desc='Instrumental covariance', unit='calculation') as progress:
    detector = DetectorFunction(
        fs=fs, orbit_file=orbit_file, orbit_interp_order=5)
    instrument_covariance, _ = detector.noise_covariance(
        frequency, t,
        acceleration_asd=N_acc,
        oms_asd=delta_x,
        generation='2nd',
        combination_type='XYZ')
    progress.update(1)
scale = np.trace(instrument_covariance, axis1=-2, axis2=-1).real/3
instrument_covariance += covariance_regularization*scale[:, None, None]*np.eye(3)
active_chunks = []

\end{lstlisting}

\linesref{139}{150} obtain the instrumental XYZ covariance from
\code{DetectorFunction}. The call requests second-generation TDI,
XYZ combinations, the supplied acceleration and optical-metrology noise
amplitudes, and fifth-order orbit interpolation. Its first return value
has shape \((K,3,3)\); the second return value, signal-filter power,
is ignored by this caller.

The covariance helper evaluates the transfer functions at
\(t=T/2\), one midpoint epoch. The source waveforms span the moving
orbit over the full observation, while this instrumental noise weighting
is a frequency-dependent matrix evaluated at that epoch. The code does
not average the instrumental covariance over all observation times.

\linesref{151}{152} add a diagonal term
\[
\mathbf S_{{\rm inst},k}\leftarrow
\mathbf S_{{\rm inst},k}
+10^{-12}\frac{\operatorname{Re}\operatorname{tr}
(\mathbf S_{{\rm inst},k})}{3}\mathbf I_3.
\]
This small regularization scales with the mean instrumental channel
power at each frequency. It is added before computing absolute SNRs and
is therefore also present in later effective covariances.
\code{active_chunks} will collect IDs that initially exceed the threshold.

\Needspace{6\baselineskip}
% SOURCE eric_snr.py 155 180
\begin{lstlisting}[language=Python,numbers=left,showlines=true,firstnumber=155]
# Save settings as datasets; no HDF5 attributes.
with h5py.File(output_file, 'x') as output:
    settings = output.create_group('settings')
    settings.create_dataset('source_catalogue', data=str(LDC_file.resolve()))
    settings.create_dataset('observation_duration_s', data=T_obs)
    settings.create_dataset('detector_sampling_rate_hz', data=fs)
    settings.create_dataset('frequency_spacing_hz', data=1/T_obs)
    settings.create_dataset('smoothing_half_width_hz', data=128/T_obs)
    settings.create_dataset('signal_time_window', data='rectangular full observation')
    settings.create_dataset('signal_window_mean_square', data=1.)
    settings.create_dataset('tdi_channels', data='X2,Y2,Z2')
    settings.create_dataset('snr_threshold', data=SNR_threshold)
    settings.create_dataset('transfer_factor_threshold', data=1)
    settings.create_dataset('neighboring_bins', data=128)
    settings.create_dataset('noise_acceleration_asd', data=N_acc)
    settings.create_dataset('noise_oms_asd', data=delta_x)
    settings.create_dataset('covariance_regularization', data=covariance_regularization)
    settings.create_dataset('waveform_device', data='gpu')
    settings.create_dataset('post_waveform_device', data='gpu')
    settings.create_dataset('snr_cpu_threads', data=0)

    status = output.create_group('status')
    status.create_dataset('complete', data=False)
    status.create_dataset('source_count', data=source_count)
    output.create_group('sky')

\end{lstlisting}

\linesref{155}{179} create the temporary working HDF5 file and record its
settings as datasets. Mode \code{'x'} creates a new working file in the
fresh temporary directory. \code{settings} stores the source path,
duration, sampling, Fourier spacing, smoothing width, rectangular window,
channels, threshold, noise amplitudes and regularization.

The device entries describe the waveform and subsequent numerical stages.
\code{snr_cpu_threads=0} records that the SNR work is dispatched to the
GPU; it does not mean that host file reading and bookkeeping use no CPU.
\code{transfer_factor_threshold=1} records an integer stopping convention:
the iteration ends when fewer than one source transfers, which the
loop implements explicitly as \code{transfer_factor == 0}.

\code{status/complete} begins as false and \code{status/source_count}
records the total number of rows. The \code{sky} group holds the original
catalogue-layout working datasets created in the next block. These
compatibility datasets coexist with the final single numerical table;
the current subtraction reader uses \code{catalogue}.

\Needspace{6\baselineskip}
% SOURCE eric_snr.py 181 210
\begin{lstlisting}[language=Python,numbers=left,showlines=true,firstnumber=181]
fields = ['Name', *PARAMETER_FIELDS]
number_chunks = sum((size+catalogue_chunk_size-1)//catalogue_chunk_size
                    for size in catalogue_sizes.values())
gpu = GPUJobs(T_obs, fs, orbit_file, first_bin)
chunk_index = 0
with tqdm(total=source_count, desc='Round 1: all absolute SNRs',
          unit='DWD', unit_scale=True) as progress:
    for group in catalogue_groups:
        dataset_path = f'sky/{group}/cat'
        group_size = catalogue_sizes[group]
        with h5py.File(LDC_file, 'r') as ldc:
            catalogue_dtype = ldc[dataset_path].dtype
        # Preserve source IDs; store every numerical field without decimal rounding.
        dtype = np.dtype([('Name', catalogue_dtype.fields['Name'][0])]
                         + [(name, '<f8') for name in PARAMETER_FIELDS]
                         + [('SNR_absolute', '<f8'), ('SNR_relative', '<f8')])
        chunk_shape = (min(4096, group_size), 1)
        with h5py.File(output_file, 'a') as output:
            catalogue_group = output['sky'].create_group(group)
            catalogue_group.create_dataset('cat', shape=(group_size, 1), dtype=dtype,
                                           chunks=chunk_shape, compression='lzf', shuffle=True)

        for first in range(0, group_size, catalogue_chunk_size):
            chunk_index += 1
            last = min(first+catalogue_chunk_size, group_size)
            # Read one chunk.
            with h5py.File(LDC_file, 'r') as ldc:
                records = ldc[dataset_path].fields(fields)[first:last, 0]

            # Calculate with the files closed.
\end{lstlisting}

\linesref{181}{185} choose the original \code{Name} field and eight
physical parameters, count the outer processing chunks, and instantiate
one \code{GPUJobs} object. With 26 million and 3 million input rows,
the 65,536-row setting produces 397 and 46 chunks, respectively:
443 chunks in total. This is a file-processing count, not a count
of observation segments or simultaneous GPU sources.

\linesref{188}{201} preserve each input dataset's original \code{Name}
dtype and store the eight parameters and two SNRs as float64 fields.
Each compatibility dataset has shape \((n_{\rm group},1)\) and a compound
dtype; its named fields are not a numerical second dimension.
This explains why the separate numerical \((29{,}000{,}000,11)\)
table is created in Step 4.

\code{chunk_shape} controls HDF5 storage chunks of at most 4,096 rows.
LZF compression and byte shuffling are lossless. The outer
\code{catalogue_chunk_size} separately controls the 65,536-row reads.
Lines 203--208 load one chunk's requested fields. The loading scope closes
before the waveform and SNR calculations begin in the next block.

\Needspace{6\baselineskip}
% SOURCE eric_snr.py 211 233
\begin{lstlisting}[language=Python,numbers=left,showlines=true,firstnumber=211]
            rows = np.empty(len(records), dtype=dtype)
            for name in fields:
                rows[name] = records[name]
            rows['SNR_relative'] = np.nan
            dwd_ids = catalogue_offsets[group]+np.arange(first, last)
            for batch_ids, snr in gpu.foreground(
                    dwd_ids, instrument_covariance, chunk=(chunk_index, number_chunks)):
                rows['SNR_absolute'][batch_ids-catalogue_offsets[group]-first] = snr
                progress.update(len(batch_ids))
            keep = rows['SNR_absolute'] > SNR_threshold
            active_chunks.append(dwd_ids[keep])

            # Save the calculated chunk.
            with h5py.File(output_file, 'a') as output:
                output[dataset_path][first:last, 0] = rows

active = np.concatenate(active_chunks) if active_chunks else np.empty(0, dtype=np.int64)
del active_chunks
resolved_count = len(active)
history.append(dict(round=1, evaluated=source_count, resolved=resolved_count,
                    unresolved=source_count-resolved_count, transfer_factor=None))
print(f'Round 1: {resolved_count:,} resolved; {source_count-resolved_count:,} initially unresolved', flush=True)

\end{lstlisting}

\linesref{211}{215} copy the input parameters and original names into the
working rows. \code{SNR_relative} is temporarily \code{NaN}, because the
foreground model has not yet been fixed. \code{dwd_ids} contains
zero-based global indices, including the second dataset's offset.

\linesref{216}{221} call \code{gpu.foreground} with the instrumental
covariance alone. The returned values are saved as \code{SNR_absolute};
the generator's global batch indices are converted back to positions in
the current row chunk. At the same time, the GPU helper accumulates
sources with absolute SNR at or below seven into the initial foreground.
The rows with absolute SNR strictly above seven supply the initial
\code{active} IDs.

\linesref{223}{225} write the calculated rows with the source file closed.
Lines 227--232 join the active-ID chunks and record round 1.
Its \code{evaluated} count is the full source count, \code{resolved}
is the initial active count and \code{unresolved} is the complement.
\code{transfer_factor=None} marks that there is no earlier resolved set
to transfer from in this first pass.

In this producer, \emph{resolved} means above the model's SNR threshold.
It does not establish that RJMCMC has identified or fitted those sources.
All original parameter rows are retained in the output, regardless of
their threshold classification.

\Needspace{6\baselineskip}
% SOURCE eric_snr.py 234 254
\begin{lstlisting}[language=Python,numbers=left,showlines=true,firstnumber=234]
instrument_covariance = gpu.noise_covariance
confusion_covariance = gpu.confusion_covariance()
with h5py.File(background_file, 'x') as background:
    settings = background.create_group('settings')
    settings.create_dataset('observation_duration_s', data=T_obs)
    settings.create_dataset('detector_sampling_rate_hz', data=fs)
    settings.create_dataset('frequency_spacing_hz', data=1/T_obs)
    settings.create_dataset('smoothing_half_width_hz', data=128/T_obs)
    settings.create_dataset('signal_time_window', data='rectangular full observation')
    settings.create_dataset('signal_window_mean_square', data=1.)
    settings.create_dataset('covariance_axes', data='frequency,XYZ,XYZ')
    settings.create_dataset('neighboring_bins', data=128)

    status = background.create_group('status')
    status.create_dataset('complete', data=False)
    background.create_dataset('frequency_hz', data=frequency)
    background.create_dataset('instrument_covariance', data=np.asarray(instrument_covariance), compression='lzf')
    background.create_dataset('initial_confusion_covariance',
                              data=np.asarray(confusion_covariance),
                              compression='lzf')

\end{lstlisting}

\linesref{234}{235} retain the instrumental covariance already copied
to the GPU and smooth the initial unresolved spectral accumulator.
\linesref{236}{253} save a diagnostic \code{background.h5} inside the
temporary directory. Its frequency array has shape \((K,)\); the
instrumental and initial confusion arrays have shape \((K,3,3)\).
\code{np.asarray} copies them to the host for HDF5 storage.

The background settings describe the frequency grid, window and channel
axes. Its completion flag begins as false, and its final covariances
and round history are added at the end. The script subsequently deletes
this temporary background file when it removes the working directory.
The published \code{ldc_dwd_data.h5} is the SNR catalogue and its
settings, status and round history.

\Needspace{18\baselineskip}
\subsection{Step 3: update the unresolved foreground (lines 256--283)}
\Needspace{6\baselineskip}
% SOURCE eric_snr.py 255 282
\begin{lstlisting}[language=Python,numbers=left,showlines=true,firstnumber=255]
# -----------------------------------------------------------------------
# 3. Re-evaluate only resolved sources and accumulate each transfer once.
# -----------------------------------------------------------------------
print('3/4  Iteratively update the unresolved foreground', flush=True)
round_number = 2
while True:
    effective_covariance = instrument_covariance+confusion_covariance
    evaluated = resolved_count
    transfer_factor = 0
    remaining = []
    with tqdm(total=evaluated, desc=f'Round {round_number}: resolved SNRs',
              unit='DWD', unit_scale=True) as progress:
        for batch_ids, snr in gpu.foreground(active, effective_covariance, chunk=(1, 1)):
            transfer = snr <= SNR_threshold
            transfer_factor += np.count_nonzero(transfer)
            remaining.append(batch_ids[~transfer])
            progress.update(len(batch_ids))
    active = np.concatenate(remaining) if remaining else np.empty(0, dtype=active.dtype)
    resolved_count = evaluated-transfer_factor
    history.append(dict(round=round_number, evaluated=evaluated, resolved=resolved_count,
                        unresolved=source_count-resolved_count, transfer_factor=transfer_factor))
    print(f'Round {round_number}: transfer_factor={transfer_factor:,}; '
          f'{resolved_count:,} resolved; {source_count-resolved_count:,} unresolved', flush=True)
    if transfer_factor == 0:
        break
    confusion_covariance = gpu.confusion_covariance()
    round_number += 1

\end{lstlisting}

\linesref{259}{281} iterate from round 2 with
\[
\mathbf S_{\rm eff}^{(r)}(f)
=\mathbf S_{\rm inst}(f)+\mathbf C_{\rm unresolved}^{(r-1)}(f).
\]
Only the currently active sources are re-evaluated. They were above
threshold under the previous covariance; adding foreground noise can
reduce their SNRs. A source with relative SNR at or below seven transfers
to the unresolved set. \code{transfer_factor} is the number of such
sources, and \code{remaining} collects IDs still strictly above seven.

The same \code{gpu.foreground} call adds the transferring sources'
spectral matrices to the persistent accumulator. Already unresolved
sources are absent from \code{active}, so their contribution is not added
again. The active list can only shrink; this loop has no transition
that moves an unresolved source back into it.

Lines 273--277 record the counts for that round. Here \code{evaluated}
is the previous resolved count, and the new resolved count is
\code{evaluated - transfer_factor}. The total unresolved count is
\code{source_count - resolved_count}.

If no source transfers, lines 278--279 stop the iteration. The foreground
accumulator has not changed, so the \code{effective_covariance} calculated
at the start of that round is already the final weighting. Otherwise,
line 280 smooths the updated accumulator and the next round begins.
There is no fixed maximum number of rounds in this script.

\Needspace{18\baselineskip}
\subsection{Step 4: final relative SNRs and the 11-column table (lines 284--321)}
\Needspace{6\baselineskip}
% SOURCE eric_snr.py 283 294
\begin{lstlisting}[language=Python,numbers=left,showlines=true,firstnumber=283]
# -----------------------------------------------------------------------
# 4. With the final foreground fixed, calculate every relative SNR again.
# -----------------------------------------------------------------------
print('4/4  Final relative SNRs for every LDC DWD', flush=True)
catalogue_columns = ('ID', *PARAMETER_FIELDS, 'SNR_absolute', 'SNR_relative')
with h5py.File(output_file, 'a') as output:
    output.create_dataset('catalogue', shape=(source_count, 11), dtype='f8',
                          chunks=(min(4096, source_count), 11), compression='lzf', shuffle=True)
    output['settings'].create_dataset('catalogue_columns',
        data=np.asarray(catalogue_columns, dtype=h5py.string_dtype('utf-8')))
# No transfers occurred in the last round, so its covariance is already final.
chunk_index = 0
\end{lstlisting}

\linesref{287}{292} create the public numerical \code{catalogue} dataset.
Its shape is \((\code{source_count},11)\), currently
\((29{,}000{,}000,11)\), with float64 values. Storage chunks contain at
most 4,096 complete rows; LZF and shuffling remain lossless.
\code{settings/catalogue_columns} saves the column order as UTF-8 strings.

{\small
\begin{tabularx}{\textwidth}{P{0.09\textwidth}P{0.31\textwidth}X}
\toprule Column & Saved name & Meaning\\\midrule
0 & \code{ID} & One-based global catalogue row ID.\\
1 & \code{Amplitude} & Dimensionless gravitational-wave amplitude.\\
2 & \code{Frequency} & Reference frequency in Hz.\\
3 & \code{FrequencyDerivative} & Frequency derivative in Hz/s.\\
4 & \code{InitialPhase} & Initial phase in radians.\\
5 & \code{Inclination} & Orbital inclination in radians.\\
6 & \code{EclipticLongitude} & Ecliptic longitude in radians.\\
7 & \code{EclipticLatitude} & Ecliptic latitude in radians.\\
8 & \code{Polarization} & Polarization angle in radians.\\
9 & \code{SNR_absolute} & SNR weighted by instrumental covariance.\\
10 & \code{SNR_relative} & SNR weighted by final instrumental plus foreground covariance.\\
\bottomrule
\end{tabularx}}

Each ID is represented exactly in float64 at this catalogue size.
The raw numerical table occupies \(29{,}000{,}000\times11\times8
=2{,}552{,}000{,}000\) bytes before compression.
The next pass fills all its rows; this dataset is not restricted to
above-threshold DWDs.

\Needspace{6\baselineskip}
% SOURCE eric_snr.py 295 320
\begin{lstlisting}[language=Python,numbers=left,showlines=true,firstnumber=295]
with tqdm(total=source_count, desc='Final round: all relative SNRs',
          unit='DWD', unit_scale=True) as progress:
    for group in catalogue_groups:
        dataset_path = f'sky/{group}/cat'
        group_size = catalogue_sizes[group]
        for first in range(0, group_size, catalogue_chunk_size):
            chunk_index += 1
            last = min(first+catalogue_chunk_size, group_size)
            # Read the previously saved absolute SNRs.
            with h5py.File(output_file, 'r') as output:
                rows = output[dataset_path][first:last, 0]

            # Calculate with the files closed.
            dwd_ids = catalogue_offsets[group]+np.arange(first, last)
            for batch_ids, snr in gpu.snr(
                    dwd_ids, effective_covariance, chunk=(chunk_index, number_chunks)):
                rows['SNR_relative'][batch_ids-catalogue_offsets[group]-first] = snr
                progress.update(len(batch_ids))

            merged_rows = np.column_stack((dwd_ids + 1,
                *[rows[name] for name in catalogue_columns[1:]]))
            # Save the calculated chunk.
            with h5py.File(output_file, 'a') as output:
                output[dataset_path][first:last, 0] = rows
                output['catalogue'][catalogue_offsets[group] + first:catalogue_offsets[group] + last] = merged_rows

\end{lstlisting}

\linesref{295}{319} re-evaluate every source with the final
\code{effective_covariance}. This full pass is required because
initially unresolved sources and sources transferred in earlier rounds
have not all been evaluated against that same final covariance.

Lines 303--305 load the previously saved parameters, original name
and absolute SNR. \code{gpu.snr}, rather than \code{gpu.foreground},
then calculates the final \code{SNR_relative} without adding to the
foreground. The absolute SNR remains the Step 2 result.

Lines 314--315 assemble each numerical row as its global index plus one,
the eight parameters, and the two SNRs. This conversion changes the
table layout, not the saved parameter or SNR values. Lines 317--319
write both the compatibility rows and the corresponding slice of
the single merged table. The original \code{Name} values remain in
the compatibility datasets, while \code{ID} in the merged table follows
global row order.

The subtraction pipeline reads this table directly and retains its
rows with column 10 strictly above \code{SNR_r_threshold}; see
Section~\ref{sec:observation_preparation}. It does not need to handle
the original storage-group names or rebuild IDs. Both SNRs refer to
the same full-duration source response, with different covariance
weightings as defined in Equation~\ref{eq:snr_producer_normalization}.

\Needspace{30\baselineskip}
\subsection{Record the rounds, publish the result and clean up (lines 321--360)}
\Needspace{26\baselineskip}
% SOURCE eric_snr.py 321 343
\begin{lstlisting}[language=Python,numbers=left,showlines=true,firstnumber=321]
# Store one value per round in each dataset.
round_numbers = np.array([row['round'] for row in history])
evaluated_counts = np.array([row['evaluated'] for row in history])
resolved_counts = np.array([row['resolved'] for row in history])
unresolved_counts = np.array([row['unresolved'] for row in history])
# Round 1 has no transfers yet; NaN means not applicable.
transfer_counts = np.array([
    np.nan if row['transfer_factor'] is None else row['transfer_factor'] for row in history])
runtime_seconds = time.perf_counter()-started

with h5py.File(output_file, 'a') as output:
    rounds = output.create_group('rounds')
    rounds.create_dataset('round', data=round_numbers)
    rounds.create_dataset('evaluated', data=evaluated_counts)
    rounds.create_dataset('resolved', data=resolved_counts)
    rounds.create_dataset('unresolved', data=unresolved_counts)
    rounds.create_dataset('transfer_factor', data=transfer_counts)

    status = output['status']
    status.create_dataset('resolved_count', data=resolved_count)
    status.create_dataset('runtime_seconds', data=runtime_seconds)
    status['complete'][()] = True

\end{lstlisting}

\linesref{321}{329} convert the in-memory round history into arrays.
There is one entry per round for its number, evaluated count,
resolved count, unresolved count and transfer count.
The first transfer entry is saved as \code{NaN}, matching its
\emph{not applicable} meaning in Step 2. The runtime is measured from
the producer's initial \code{started} value.

\linesref{331}{342} save these arrays under \code{rounds}, add the final
resolved count and elapsed seconds under \code{status}, and mark the
working catalogue complete. At this point all 29 million rows in
\code{catalogue} have their absolute and final relative SNRs.
The public output has no per-source waveform or strain dataset;
the large individual XYZ batches were used to calculate these scalars.

\Needspace{6\baselineskip}
% SOURCE eric_snr.py 344 360
\begin{lstlisting}[language=Python,numbers=left,showlines=true,firstnumber=344]
with h5py.File(background_file, 'a') as background:
    background.create_dataset('final_confusion_covariance', data=np.asarray(confusion_covariance), compression='lzf')
    background.create_dataset('final_effective_covariance', data=np.asarray(effective_covariance), compression='lzf')

    rounds = background.create_group('rounds')
    rounds.create_dataset('round', data=round_numbers)
    rounds.create_dataset('evaluated', data=evaluated_counts)
    rounds.create_dataset('resolved', data=resolved_counts)
    rounds.create_dataset('unresolved', data=unresolved_counts)
    rounds.create_dataset('transfer_factor', data=transfer_counts)
    background['status/complete'][()] = True
# Publish only the completed catalogue, then remove the temporary files.
output_file = output_file.replace(final_output_file)
temporary_directory.cleanup()
print(f'Saved: {output_file.resolve()}', flush=True)
print(f'{resolved_count:,}/{source_count:,} above SNR_r > {SNR_threshold:g}; '
      f'{time.perf_counter()-started:.1f} s', flush=True)
\end{lstlisting}

\linesref{344}{355} complete the temporary background diagnostics with
the final confusion covariance, effective covariance and round counts.
\linesref{356}{360} then move the completed working catalogue to
\code{final_output_file}. \code{Path.replace} replaces an existing result
at that destination. Because the temporary directory was created under
the output parent, the working and final files use the same filesystem.

After publication, \code{temporary_directory.cleanup()} removes the
temporary directory and \code{background.h5}. The final messages display
the saved catalogue path, resolved count, threshold and elapsed time.
The durable result contains \code{catalogue}, column names and settings,
completion status, round history and compatibility catalogue datasets.
The numerical table is the interface used by the current subtraction code.

This review reproduces the inspected local and server source, with SHA-256
\path{7441c8ba7b7e0dc4680e5e4a14b6bb435be07684c68816f300911d46f1b0510e}.
The table export and direct reader were checked in the recorded regression
tests in Section~\ref{sec:verification}. Adding this review changes the
documentation; it does not rerun the 29-million-source SNR calculation.



\clearpage
\Needspace{8\baselineskip}
\section{Subtraction pipeline: imports and configuration}
\label{sec:walkthrough1}

The primary reviewed file is \mainfile. The pipeline supports the full
frequency grid and a restricted subtraction interval. The three input files are the observation's companion
orbit, the completed \code{eric_snr.py} result and the simulated observation.
Section~\ref{sec:simulation_pipeline} describes the producer of that observation.
Section~\ref{sec:snr_pipeline} reviews the producer of the SNR catalogue.

The main subtraction pipeline is documented as Steps 1--11, with every source
line retained. Step numbers in Section~\ref{sec:simulation_pipeline} refer to
the separate observation simulation.


% SOURCE eric_subtraction_all.py 1 24
\begin{lstlisting}[language=Python,numbers=left,firstnumber=1,showlines=true]
'''Step 1: import all the necessary packages.'''
from pathlib import Path
from tempfile import mkdtemp
import hashlib
import json
import shutil
import h5py
import numpy as np
import lisagwresponse
import lisainstrument
import pytdi
from eric_generate_waveform import DWD_simulation_gpu, tdi_gpu
from script.eric_Fstats import (
    build_response, physical_amplitude_coordinates,
    physical_amplitudes, reconstruct_waveform, noise_inner_product,
)
from script.dwd_blind_search import (
    residual_covariance, template_f_statistic_search, amplitude_f_statistic,
)
from script.dwd_bank_association import assign_bank_ids
from eric_rjmcmc import run_eryn_rjmcmc
from script.oneblock_observation import read_checkpoint, write_checkpoint
from script.tool import _save_tree, single_cpu

\end{lstlisting}

\textbf{Step 1.} Import all the packages required by the subtraction workflow.


% SOURCE eric_subtraction_all.py 25 38
\begin{lstlisting}[language=Python,firstnumber=25,numbers=left,numberblanklines=true,showlines=true]
'''Step 2: define the parameters.'''
orbit_file = Path('./data/dwd_observation_full_catalogue.orbits.h5')
snr_file = Path('./data/ldc_dwd_data.h5')
observation_file = Path('./data/dwd_observation_full_catalogue.h5')
output_file = Path('./data/dwd_subtraction_all.h5')
frequency_range = (0.0001, 0.05)
f_derivative_pn = [-25, -10]
freq_block_width = 1e-05
subtraction_frequency_range = None
frequency_derivative_bounds = 10. ** f_derivative_pn[1] * np.array([-1., 1.])
F_threshold = 16.0
SNR_r_threshold = 7.0
sample_number = (500, 1500)
maximum_sources = 64
\end{lstlisting}

\textbf{Step 2.} Define the input paths, search limits,
detection thresholds and sampling effort.

\Needspace{6\baselineskip}



The defaults below are read from this source snapshot. Sampling effort and a
source-capacity setting specify computational work; neither establishes convergence.

{\small
\begin{longtable}{P{0.31\textwidth}P{0.22\textwidth}P{0.39\textwidth}}
\toprule Setting & Default & Meaning\\\midrule
\endfirsthead
\toprule Setting & Default & Meaning\\\midrule
\endhead
\code{T_obs}, \code{fs} & From observation file & Saved duration and sampling frequency.\\
\code{freq_block_width} & \(10^{-5}\Hz\) & Width of the fixed carrier-frequency blocks.\\
\code{frequency_range} & \([10^{-4},0.05]\) Hz & Limits of the indexed frequency grid.\\
\code{f_derivative_pn} & \code{[-25, -10]} & Declared common \(\log_{10}|\dot f|\) interval for positive and negative derivatives, using values in Hz/s.\\
\code{subtraction_frequency_range} & \code{None} & Requested interval; \code{None} selects all grid blocks.\\
\code{F_threshold} & 16 & Threshold on \(\mathcal F\), rather than \(2\mathcal F\).\\
\code{SNR_r_threshold} & 7 & Saved relative SNR for candidate selection and fitted-source SNR must be strictly larger.\\
\code{sample_number} & \code{(500, 1500)} & Burn-in and production cycles for each sampler call.\\
\code{maximum_sources} & 64 & Minimum per-fit sampler capacity.\\
\bottomrule
\end{longtable}}


% SOURCE eric_subtraction_all.py 39 56
\begin{lstlisting}[language=Python,firstnumber=39,numbers=left,showlines=true]
with h5py.File(observation_file, 'r') as saved:
    T_obs = saved['observation'].attrs['duration']
    fs = saved['observation'].attrs['rate']
    total_xyz_td = saved['observation/total_xyz_td'][:]
    detector_config = json.loads(saved['observation'].attrs['detector_config'])
detector_config['orbits'] = str(orbit_file)
N, N_day = int(round(T_obs * fs)), int(round(86400 * fs))
frequency_blocks = np.arange(
    frequency_range[0], frequency_range[1] - freq_block_width * 1.e-10, freq_block_width,
)[:, None] + np.array([0., freq_block_width])
frequency_blocks[-1, 1] = frequency_range[1]
block_count = len(frequency_blocks)
block_indices = np.arange(block_count, dtype=np.int64)
subtraction_range = np.array([frequency_range[0], frequency_range[1]]) if subtraction_frequency_range is None else np.asarray(subtraction_frequency_range, dtype=float)
selected_block_indices = block_indices[
    (frequency_blocks[:, 0] >= subtraction_range[0] - freq_block_width * 1.e-10)
    & (frequency_blocks[:, 1] <= subtraction_range[1] + freq_block_width * 1.e-10)]
selected_block_count = len(selected_block_indices)
\end{lstlisting}

Lines 39--56 read the observation and its saved duration and sampling rate, then construct the indexed grid and select the requested blocks.

The grid is built once, before any global round. By default its lower edge is
\(10^{-4}\) Hz, its upper edge is \(0.05\) Hz, and the last block may be shorter
than the configured width. At the default block width it contains 4,990
blocks with indices \(0,\ldots,4989\). Changing the subtraction interval leaves
these edges and indices unchanged.

The subtraction interval is chosen to cover complete blocks. For a requested
interval \([a,b]\), block \(j\) is selected when
\[
f_j^{\rm lo}\ge a,\qquad f_j^{\rm hi}\le b.
\]
The small numerical tolerance only handles floating-point boundary rounding.
The selected indices refer to the original \code{frequency_blocks} rows.
Template evaluation and RJMCMC use each complete block's original bounds;
there are no clipped intervals or separate \code{search_range} variables.
The final grid upper edge is included once.

The caller creates the output's parent directory.



% SOURCE eric_subtraction_all.py 57 57
\begin{lstlisting}[language=Python,firstnumber=57,numbers=left,showlines=true]
fields = dict(A=1, f=2, df=3, phi0=4, iota=5, longitude=6, latitude=7, psi_bme=8)
\end{lstlisting}

Line 57 maps each physical parameter to its column position in the numerical SNR table.

The dictionary provides the eight physical parameter column positions used for
recovered waveforms and post-fit labels. Column 0 stores the ID, columns 1--8
store the physical parameters, and columns 9--10 store the absolute and relative SNRs. All DWDs use the same waveform and inference
procedure. The detector configuration is read directly from the observation; its filter,
clock, padding and GPU settings are reused for fitted-waveform synthesis.
Only the orbit-file location follows the separately specified companion path.

{\small
\begin{tabularx}{\textwidth}{P{0.14\textwidth}P{0.08\textwidth}P{0.26\textwidth}X}
\toprule Parameter & Symbol & Catalogue field & Description\\\midrule
\code{A} & \(A\) & \code{Amplitude} & Gravitational-wave strain amplitude.\\
\code{f} & \(f\) & \code{Frequency} & Reference gravitational-wave frequency.\\
\code{df} & \(\dot f\) & \code{FrequencyDerivative} & Frequency derivative.\\
\code{phi0} & \(\phi_0\) & \code{InitialPhase} & Initial waveform phase.\\
\code{iota} & \(\iota\) & \code{Inclination} & Orbital inclination to the line of sight.\\
\code{longitude} & \(\lambda\) & \code{EclipticLongitude} & Ecliptic sky longitude.\\
\code{latitude} & \(\beta\) & \code{EclipticLatitude} & Ecliptic sky latitude.\\
\code{psi_bme} & \(\psi\) & \code{Polarization} & Polarization angle in the barycentric mean ecliptic frame.\\
\bottomrule
\end{tabularx}}


\Needspace{8\baselineskip}
\section{Shared input, observation and result preparation}
\label{sec:observation_preparation}

\Needspace{8\baselineskip}
\subsection{Step 3: read the inputs and prepare shared state (lines 58--73)}


% SOURCE eric_subtraction_all.py 58 73
\begin{lstlisting}[language=Python,firstnumber=58,numbers=left,showlines=true]
# Step 3. read the candidate bank and prepare shared state.
with h5py.File(snr_file, 'r') as saved:
    template_dwd = saved['catalogue'][saved['catalogue'][:, 10] > SNR_r_threshold]
total_xyz_fd, fft_frequency = np.fft.rfft(total_xyz_td, axis=-1), np.fft.rfftfreq(N, d=1 / fs)
work = Path(mkdtemp(prefix='.dwd_all_', dir=output_file.parent)).resolve()
configuration = dict(duration=T_obs, rate=fs, frequency_range=np.array([frequency_range[0], frequency_range[1]]),
    frequency_blocks_hz=frequency_blocks, block_indices=block_indices, frequency_block_width=freq_block_width,
    subtraction_frequency_range=subtraction_range, selected_block_indices=selected_block_indices,
    derivative_bounds=frequency_derivative_bounds,
    chirp_coordinate='linear', f_threshold=F_threshold,
    relative_snr_threshold=SNR_r_threshold, sample_number=np.asarray(sample_number),
    minimum_sampler_capacity=maximum_sources)
with h5py.File(output_file, 'w') as saved:
    _save_tree(saved.create_group('config'), configuration)
    saved.create_group('blocks')
    saved.create_group('rounds')
\end{lstlisting}

The result of \code{eric_snr.py} contains one numerical HDF5 dataset,
\code{catalogue}, with shape \((29{,}000{,}000,11)\). Its columns are
\code{ID}, \code{Amplitude}, \code{Frequency}, \code{FrequencyDerivative},
\code{InitialPhase}, \code{Inclination}, \code{EclipticLongitude},
\code{EclipticLatitude}, \code{Polarization}, \code{SNR_absolute}, and
\code{SNR_relative}, in that order. Column names are also saved in
\code{settings/catalogue_columns}. All values use float64; the integer IDs
1--29,000,000 are represented exactly.

The loading block reads column 10 and directly selects complete rows with
\mbox{\code{SNR_relative}} \(>\) \mbox{\code{SNR_r_threshold}}, currently 7.
A relative SNR equal to 7 is excluded. The retained \code{template_dwd} has
shape \((n_{\rm template},11)\); the current result supplies 8,336 rows.
The saved IDs are retained without renumbering or reconstructing them.
The main pipeline passes selected rows of this same numerical table directly
to the F-statistic and RJMCMC functions.
Those functions read the required intrinsic columns internally without
changing the input table. Templates are selected by the requested block
frequency limits. The physical refinement bounds now come from the upper
log-magnitude limit in \code{f_derivative_pn}: currently
\([-10^{-10},10^{-10}]\) Hz/s. The sampler retains its signed linear
derivative coordinate; the lower declared magnitude \(10^{-25}\) is not
a cutoff in this linear prior. These bounds do not generate additional templates.

The observation must contain finite \code{observation/total_xyz_td} with shape
\((3,N)\), where \(N=\mathrm{round}(T_{\rm obs}f_s)\). The duration and sampling rate are read from \code{observation} attributes
\code{duration} and \code{rate}; the subtraction script defines neither again.
The original total observation is kept unchanged. Its unscaled real FFT,
\code{total_xyz_fd}, and matching frequency axis are assigned together here.
NumPy supplies the spectrum through \code{rfft} and the frequencies through
\code{rfftfreq}. The frequency axis is reused throughout the frequency-block
analysis.

The temporary directory \code{.dwd_all_*} holds response wrappers, post-fit
label caches and sampler checkpoints. It is removed after successful
finalization. The result is opened with mode \code{'w'}, so a rerun replaces an
existing result. Configuration records the whole indexed grid as well as
\code{selected_block_indices}; the selected subset is therefore recoverable from
the result without guessing from populated block groups.

\Needspace{8\baselineskip}
\section{Global rounds and shared residual covariance}
\label{sec:architecture}

Steps 4--10 repeat for each global round. Steps 5--8 visit each active selected
block; Steps 9--10 combine the accepted sources and decide whether another
round is needed.

\Needspace{8\baselineskip}
\subsection{Step 4: start a round and estimate covariance (lines 74--95)}


% SOURCE eric_subtraction_all.py 74 95
\begin{lstlisting}[language=Python,firstnumber=74,numbers=left,showlines=true]
# Step 4. start a global round and estimate covariance.
round_number, last_accepted_round = 0, 0
print(f'Subtraction range {subtraction_range[0]:.8g}--{subtraction_range[1]:.8g} Hz: '
      f'selected {selected_block_count} of {block_count} blocks, indices '
      f'{selected_block_indices[0]}--{selected_block_indices[-1]} (zero-based).', flush=True)
with single_cpu():
    while True:
        round_number += 1
        label = f'{round_number:02d}'
        print(f'Global round {round_number}: visit {selected_block_count} selected frequency blocks.', flush=True)
        # This complete Fourier snapshot is unchanged throughout the block loop.
        round_residual_fd = total_xyz_fd if round_number == 1 else residual_xyz_fd
        round_residual_fd.setflags(write=False)
        round_digest = hashlib.sha256(round_residual_fd.tobytes()).hexdigest()
        covariance_full = residual_covariance(round_residual_fd, T_obs, fs)
        block_order = selected_block_indices
        status = np.full(block_count, 'outside_subtraction_range', dtype='S32')
        status[selected_block_indices] = 'pending'
        with h5py.File(output_file, 'a') as saved:
            group = saved['rounds'].create_group(label)
            group.create_group('blocks')
            group.attrs.update(complete=False, committed=False, residual_input_sha256=round_digest)
\end{lstlisting}

The first round analyzes \code{total_xyz_fd}. Later rounds analyze the
residual left by the preceding subtraction round.


Every selected block fits the same full Fourier residual snapshot in a global
round. The snapshot is made read-only and its digest is recorded before the
block loop. All accepted candidates remain proposals until the selected pass
is complete.

The covariance helper estimates the XYZ noise covariance from overlapping
measured-data chunks. Every requested block enters the F-statistic search in
index order; there is no spectral-excess or false-alarm pre-screen. Outside
blocks retain the status \code{outside_subtraction_range}.

\Needspace{8\baselineskip}
\section{Block preparation and template-bank F-statistics}
\label{sec:bank_preparation}

\Needspace{8\baselineskip}
\subsection{Step 5: prepare each selected block and evaluate its templates (lines 96--144)}


% SOURCE eric_subtraction_all.py 96 144
\begin{lstlisting}[language=Python,firstnumber=96,numbers=left,showlines=true]
        # Step 5. evaluate each selected block against its recorded templates.
        selected_counts = np.zeros(block_count, dtype=np.int64)
        accepted_counts = np.zeros(block_count, dtype=np.int64)
        proposals, proposal_doppler, proposal_coefficients, proposal_ids, proposal_blocks = [], [], [], [], []
        for block_index in block_order:
            key = f'{block_index:05d}'
            block_range = frequency_blocks[block_index]
            center = float(block_range.mean())
            block_derivative_bounds = np.asarray(frequency_derivative_bounds, dtype=float)
            padding = block_range[1] * 1.e-4 + max(abs(block_derivative_bounds)) * T_obs + 8 / T_obs
            retained_bins = np.flatnonzero((fft_frequency >= block_range[0] - padding)
                                          & (fft_frequency <= block_range[1] + padding) & (fft_frequency > 0))
            center_bin = int(round(center * T_obs))
            wide_bins = np.arange(max(1, min(center_bin - 2048, retained_bins[0] - 256)),
                min(len(fft_frequency), max(center_bin + 2049, retained_bins[-1] + 257)))
            core_indices = retained_bins - wide_bins[0]
            physical_frequency, wide_frequency = fft_frequency[retained_bins], fft_frequency[wide_bins]
            if round_number == 1 and block_index == block_order[0]:
                builder, response_times, response_geometry = build_response(work, wide_frequency, T_obs, fs, orbit_file=orbit_file)
                builder_type = type(builder)
            else:
                builder = builder_type(np.broadcast_to(np.eye(3), (len(wide_frequency), 3, 3)),
                    np.zeros((3, len(wide_frequency)), complex), response_times, response_geometry, wide_frequency, T_obs, fs)
            template_data = dict(derivative_bounds=block_derivative_bounds,
                retained_bins=retained_bins, wide_bins=wide_bins, core_indices=core_indices,
                block_index=block_index, frequency_range=block_range,
                physical_frequency_hz=physical_frequency)
            with h5py.File(output_file, 'a') as saved:
                if key not in saved['blocks']:
                    _save_tree(saved['blocks'].create_group(key), template_data)
            covariance = covariance_full[retained_bins].copy()
            inverse = np.linalg.inv(covariance)
            covariance.setflags(write=False)
            data = round_residual_fd[:, retained_bins]
            target = (template_dwd[:, fields['f']] >= block_range[0]) & (
                (template_dwd[:, fields['f']] < block_range[1]) |
                ((block_index == block_count - 1) & (template_dwd[:, fields['f']] == block_range[1])))
            search = template_f_statistic_search(data, covariance, builder, core_indices, T_obs, fs,
                template_dwd[target], f_threshold=F_threshold)
            search['template_dwd_id'] = template_dwd[target, 0].astype(np.int64)
            selected = search['template_dwd']
            selected_counts[block_index] = len(selected)
            block_data = dict(fixed_covariance=covariance,
                screening=dict(search=search, residual_input_xyz_fd=data))
            block_attributes = dict(residual_input_sha256=round_digest,
                block_index=block_index,
                fixed_covariance_sha256=hashlib.sha256(covariance.tobytes()).hexdigest(),
                proposed_for_subtraction=False, accepted_for_subtraction=False)
            status[block_index] = 'no_new_candidates'
\end{lstlisting}

This step creates the counts and proposal collections for the current block
pass. Response geometry is obtained directly for the first selected block of
the first round and then reused.


Each selected block retains measured bins beyond its carrier
interval, with margin
\[
\Delta f=f_{\rm upper}\,10^{-4}
      +\max|\dot f|\,T_{\rm obs}+\frac{8}{T_{\rm obs}}.
\]
The wider response window has at least the nominal 4,097-bin support, clipped
at FFT boundaries or enlarged if the retained bins require it. This support
preserves sidebands and chirp power; it does not expand the allowed carrier interval.

The first selected block constructs moving-orbit response geometry. Later blocks
reuse the same response times and geometry with their own physical frequency
arrays. Covariance is copied from the measured full-band estimate and held
read-only during the block's fit.

\Needspace{8\baselineskip}
\subsection{Evaluate the templates recorded in eric\_snr}

For each block, \code{template_f_statistic_search} receives rows of the
unchanged 11-column \code{template_dwd} table whose carrier frequencies lie
inside the complete selected block. It reads the
\((f,\dot f,\lambda,\beta)\) columns internally and returns retained rows
in the same table format for the RJMCMC call. It builds four basis waveforms for every
template and analytically maximizes over the four amplitude coefficients:
\[
\mathcal F=\frac12 x^{\mathsf T}M^+x.
\]
There is no extra frequency, sky or chirp grid and no intrinsic-parameter
optimizer before RJMCMC. Every supplied template is evaluated against the same
block residual. Templates with \(\mathcal F\ge F_{\rm threshold}\) initialize
RJMCMC, which refines the retained parameters. If no template passes, the
search result is saved and the block has no sampler candidates.

\Needspace{8\baselineskip}
\section{Joint block inference and data-based acceptance}
\label{sec:walkthrough2}

\Needspace{8\baselineskip}
\subsection{Step 6: run RJMCMC and assess the fitted sources (lines 145--204)}


% SOURCE eric_subtraction_all.py 145 204
\begin{lstlisting}[language=Python,firstnumber=145,numbers=left,showlines=true]
            # Step 6. run RJMCMC and assess fitted sources.
            if len(selected):
                print(f'  Block {key}, {block_range[0]:.8f}--{block_range[1]:.8f} Hz: fit {len(selected)} candidates.', flush=True)
                checkpoints = work / f'round_{label}_block_{key}'
                posterior = run_eryn_rjmcmc(data, covariance, selected, block_range,
                    block_derivative_bounds, center, T_obs, fs, checkpoints,
                    chain_count=1, inverse_temperatures=(1., .8, .6, .4),
                    iterations=sum(sample_number), burn_in=sample_number[0],
                    maximum_sources=max(maximum_sources, len(selected)), chirp_coordinate='linear',
                    response_times=response_times, response_geometry=response_geometry, physical_frequency=physical_frequency)
                doppler, coefficients = posterior['doppler'], posterior['fit_result']['coefficients']
                new_bases = builder.build(doppler)[:, :, :, core_indices]
                new_sources = reconstruct_waveform(coefficients, new_bases)
                source_snr = np.sqrt(np.maximum(noise_inner_product(new_sources, new_sources, inverse, T_obs, fs), 0.))
                assessment_data = data.copy()
                # Fits share the round snapshot, but acceptance is conditional on
                # already proposed components to avoid subtracting a sideband twice.
                if proposal_doppler:
                    points, amplitudes = np.concatenate(proposal_doppler), np.concatenate(proposal_coefficients)
                    nearby = abs(points[:, 0] - center) <= .5 * np.ptp(block_range) + 2 * padding
                    if np.any(nearby):
                        pending = reconstruct_waveform(amplitudes[nearby],
                            builder.build(points[nearby])[:, :, :, core_indices])
                        assessment_data -= pending.sum(axis=0)
                conditional_f = np.zeros(len(doppler))
                for i in range(len(doppler)):
                    conditional_data = assessment_data - new_sources.sum(axis=0) + new_sources[i]
                    conditional_f[i] = amplitude_f_statistic(conditional_data, new_bases[i:i + 1],
                        covariance, T_obs, fs)['two_f'][0] / 2
                owned = (doppler[:, 0] >= block_range[0]) & ((doppler[:, 0] < block_range[1])
                    | ((block_range[1] == frequency_range[1]) & (doppler[:, 0] == frequency_range[1])))
                duplicate = np.zeros(len(doppler), dtype=bool)
                previous_doppler = proposal_doppler if round_number == 1 else accepted_doppler + proposal_doppler
                previous_coefficients = proposal_coefficients if round_number == 1 else accepted_coefficients + proposal_coefficients
                if previous_doppler:
                    points, amplitudes = np.concatenate(previous_doppler), np.concatenate(previous_coefficients)
                    nearby = abs(points[:, 0] - center) <= .5 * np.ptp(block_range) + 2 * padding
                    if np.any(nearby):
                        previous = reconstruct_waveform(amplitudes[nearby],
                            builder.build(points[nearby])[:, :, :, core_indices])
                        overlap = noise_inner_product(new_sources[:, None], previous[None, :], inverse, T_obs, fs)
                        norm = np.sqrt(np.maximum(noise_inner_product(new_sources, new_sources, inverse, T_obs, fs)[:, None]
                            * noise_inner_product(previous, previous, inverse, T_obs, fs)[None, :], 1.e-300))
                        duplicate |= np.any(overlap / norm > .9, axis=1)
                retained = owned & ~duplicate & np.isfinite(source_snr) & (source_snr > SNR_r_threshold) & (conditional_f >= F_threshold)
                # Deduplicate the fitted components using measured waveform overlap.
                for i in np.argsort(source_snr)[::-1]:
                    if not retained[i]:
                        continue
                    for j in np.flatnonzero(retained):
                        if j == i or source_snr[j] > source_snr[i] or (source_snr[j] == source_snr[i] and j < i):
                            continue
                        overlap = noise_inner_product(new_sources[i], new_sources[j], inverse, T_obs, fs)
                        if overlap / max(source_snr[i] * source_snr[j], 1.e-300) > .9:
                            retained[j], duplicate[j] = False, True
                retained_sum = new_sources[retained].sum(axis=0)
                improvement = float(noise_inner_product(assessment_data, assessment_data, inverse, T_obs, fs)
                    - noise_inner_product(assessment_data - retained_sum, assessment_data - retained_sum, inverse, T_obs, fs))
                accepted = np.any(retained) and improvement > 0
                status[block_index] = 'proposed' if accepted else 'fit_rejected_by_data'
\end{lstlisting}

Each active seeded block invokes Eryn once against the frozen residual and
fixed covariance. The default effort is 500 burn-in plus 1,500 production
cycles, with one walker and inverse temperatures \((1,0.8,0.6,0.4)\).
The controller explicitly requests the signed linear chirp coordinate.
The source capacity is at least the configured capacity and the seed count.

Fitting still sees the original round snapshot. Acceptance uses a separate
assessment array: nearby signals proposed earlier in this round are removed
from that array before judging a later proposal. This prevents neighboring
blocks from claiming the same measured feature twice without changing the
sampler's input mid-round.

For source \(i\), the conditional F statistic evaluates the remaining assessment
data with the other fitted components removed and component \(i\) retained.
An accepted component must have an owned carrier, finite fitted SNR strictly
above 7, conditional \(\mathcal F\geq16\), and no measured waveform duplicate.
Signed normalized waveform overlap greater than 0.9 detects duplicates against
previously accepted and already proposed sources; stronger members of an
overlapping fitted pair are retained.

The retained sum must also reduce the noise-weighted assessment residual norm:
\[
\Delta=\ip{y_{\rm assess}}{y_{\rm assess}}_C
       -\ip{y_{\rm assess}-h_{\rm retained}}
            {y_{\rm assess}-h_{\rm retained}}_C > 0.
\]
These decisions depend on measured data and fitted waveforms. They do not
require the recovered source to appear in a catalogue.

\Needspace{8\baselineskip}
\subsection{Step 7: associate catalogue labels after inference (lines 205--249)}


% SOURCE eric_subtraction_all.py 205 249
\begin{lstlisting}[language=Python,firstnumber=205,numbers=left,showlines=true]
                # Step 7. associate recovered sources with catalogue labels.
                # Catalogue reads happen after inference and have no acceptance role.
                if round_number == 1 and np.count_nonzero(selected_counts) == 1:
                    with h5py.File(snr_file, 'r') as saved:
                        truth_frequency = saved['catalogue'][:, fields['f']]
                        chirp = saved['catalogue'][:, fields['df']]
                    maximum_drift = float(np.max(abs(chirp) * T_obs + abs(truth_frequency) * 1.e-4 + 8 / T_obs))
                    truth_order = np.argsort(truth_frequency, kind='stable')
                    truth_frequency = truth_frequency[truth_order]
                # Truth is consulted only after fitting, and cached on disk for this block.
                truth_cache = work / f'truth_{key}.h5'
                if truth_cache.exists():
                    with h5py.File(truth_cache, 'r') as saved:
                        truth_ids = saved['dwd_id'][:]
                        truth_waveforms = saved['source_xyz_fd'][:]
                else:
                    truth_low, truth_high = block_range + np.array([-2304., 2304.]) / T_obs
                    first, last = np.searchsorted(truth_frequency, [truth_low - maximum_drift, truth_high + maximum_drift])
                    nearby_indices = np.sort(truth_order[first:last])
                    with h5py.File(snr_file, 'r') as saved:
                        records = saved['catalogue'][nearby_indices]
                    frequency = records[:, fields['f']]
                    end = frequency + records[:, fields['df']] * T_obs
                    margin = np.maximum(frequency, end) * 1.e-4 + 8 / T_obs
                    nearby = (np.minimum(frequency, end) - margin <= truth_high) & (np.maximum(frequency, end) + margin >= truth_low)
                    truth_rows = records[nearby]
                    truth_ids = truth_rows[:, 0].astype(np.int64)
                    truth_doppler = truth_rows[:, [fields[name] for name in ('f', 'df', 'longitude', 'latitude')]]
                    cos_iota = np.cos(truth_rows[:, fields['iota']])
                    truth_coefficients = physical_amplitude_coordinates(truth_rows[:, fields['A']] * (1 + cos_iota**2),
                        2 * truth_rows[:, fields['A']] * cos_iota, truth_rows[:, fields['psi_bme']], truth_rows[:, fields['phi0']]) / 1.e-22
                    truth_waveforms = np.empty((len(truth_ids), 3, len(retained_bins)), dtype=complex)
                    for first in range(0, len(truth_ids), 32):
                        truth_waveforms[first:first + 32] = reconstruct_waveform(truth_coefficients[first:first + 32],
                            builder.build(truth_doppler[first:first + 32])[:, :, :, core_indices])
                    with h5py.File(truth_cache, 'w') as saved:
                        saved.create_dataset('dwd_id', data=truth_ids)
                        saved.create_dataset('source_xyz_fd', data=truth_waveforms)
                truth_association = assign_bank_ids(new_sources, truth_waveforms, truth_ids, inverse, T_obs, fs, minimum_overlap=.9)
                matched = truth_association['matched_dwd_id'].copy()
                physical = physical_amplitudes(coefficients)
                parameters = np.column_stack((physical['A'], doppler[:, 0], doppler[:, 1], physical['phi0'],
                    physical['iota'], doppler[:, 2], doppler[:, 3], physical['psi_bme']))
                block_data.update(catalogue_association=dict(matched_dwd_id=matched,
                    best_overlap=truth_association['best_overlap'], scope='Post-fit labels only; unmatched fits may be subtracted.'),
\end{lstlisting}

The catalogue frequency lookup and its maximum frequency margin are calculated
here, when the first fitted block in the first round reaches Step 7. Both are
reused for later blocks and rounds.


Only after the fitted-source assessment does the controller scan the SNR
catalogue's frequencies and derivatives. It constructs a sorted index once,
then reads nearby rows directly from the same numerical \code{catalogue}
dataset for each block. The saved ID and eight physical parameter columns
supply the labels and reference waveforms.

The code names this temporary label view \code{truth_*}. Here it provides
post-fit labels only. Its waveforms do not seed the search, supply a prior,
control acceptance or stop the global loop. Local label waveforms are cached
for later rounds. A source absent from the file can still be accepted.

The accepted physical row is reconstructed from the fitted amplitude
coefficients and Doppler parameters, in the order
\((A,f,\dot f,\phi_0,\iota,\lambda,\beta,\psi)\). The four-amplitude inverse has
physical angle degeneracies, so this row represents an equivalent waveform
rather than a uniquely identifiable set of angles.

\Needspace{8\baselineskip}
\subsection{Step 8: allocate source IDs and record the fits (lines 250--292)}


% SOURCE eric_subtraction_all.py 250 292
\begin{lstlisting}[language=Python,firstnumber=250,numbers=left,showlines=true]
                    # Step 8. record the accepted fits.
                    acceptance=dict(source_snr=source_snr, conditional_f=conditional_f,
                        owned_by_block=owned, duplicate_previous_or_this_round=duplicate,
                        retained=retained, likelihood_improvement=improvement),
                    recovery=dict(doppler=doppler, coefficients=coefficients, parameters=parameters, source_xyz_fd=new_sources),
                    rjmcmc={name: value for name, value in posterior.items() if name not in ('bases', 'fit_result', 'checkpoint_file')})
                block_attributes.update(proposed_for_subtraction=accepted,
                    source_xyz_fd_scope='Frequency-domain fitting and post-fit association only.')
                if accepted:
                    retained_ids = matched[retained].copy()
                    if round_number == 1 and not np.any(accepted_counts):
                        next_unknown_id = -1
                    if round_number == 1:
                        used_labels = np.concatenate(proposal_ids) if proposal_ids else np.empty(0, dtype=np.int64)
                    else:
                        used_labels = np.concatenate(accepted_ids + proposal_ids)
                    for i in range(len(retained_ids)):
                        if retained_ids[i] < 0 or retained_ids[i] in used_labels or retained_ids[i] in retained_ids[:i]:
                            retained_ids[i] = next_unknown_id
                            next_unknown_id -= 1
                    proposals.append(parameters[retained])
                    proposal_doppler.append(doppler[retained])
                    proposal_coefficients.append(coefficients[retained])
                    proposal_ids.append(retained_ids)
                    proposal_blocks.append(np.full(np.count_nonzero(retained), block_index, dtype=np.int64))
                    accepted_counts[block_index] = np.count_nonzero(retained)
                snapshots = {}
                for path in sorted(checkpoints.glob('*.h5')):
                    with h5py.File(path, 'r') as saved:
                        snapshots[path.stem] = read_checkpoint(saved)
                with h5py.File(output_file, 'a') as saved:
                    group = saved[f'rounds/{label}/blocks'].create_group(key)
                    _save_tree(group, block_data)
                    group.attrs.update(block_attributes)
                    checkpoint_group = group.create_group('checkpoints')
                    for name, snapshot in snapshots.items():
                        write_checkpoint(checkpoint_group.create_group(name), snapshot)
                    group['rjmcmc'].attrs['checkpoint_group'] = f'/rounds/{label}/blocks/{key}/checkpoints'
            else:
                with h5py.File(output_file, 'a') as saved:
                    group = saved[f'rounds/{label}/blocks'].create_group(key)
                    _save_tree(group, block_data)
                    group.attrs.update(block_attributes)
\end{lstlisting}

The negative-ID counter starts when the first accepted fit needs label
allocation. The first round compares labels against its current proposals;
later rounds also include already accepted IDs.


Waveform association supplies an available positive row ID when a match exists.
Unmatched components, already used labels and repeated labels within a fit
receive fresh negative IDs. Global block indices are stored separately from
source IDs. Identity labels therefore cannot force duplicate subtraction.

Per-block records retain the search, fixed covariance, posterior output,
recovered FD components, source SNR, conditional F, ownership and duplicate
decisions, and the retained mask. Sampler HDF5 checkpoints are copied into the
result file before successful temporary-directory cleanup. The embedded
checkpoint group is recorded in the sampler attributes.

\Needspace{14\baselineskip}
\section{Combined time-domain subtraction and finalization}
\label{sec:accepted_td_subtraction}

\Needspace{8\baselineskip}
\subsection{Step 9: subtract the combined time-domain signal (lines 293--328)}


% SOURCE eric_subtraction_all.py 293 328
\begin{lstlisting}[language=Python,firstnumber=293,numbers=left,showlines=true]
        # Step 9. subtract the combined time-domain TDI signal.
        if round_number == 1:
            accepted_parameters, accepted_doppler, accepted_coefficients, accepted_ids, accepted_blocks = [], [], [], [], []
            accumulated_xyz_td = np.zeros_like(total_xyz_td)
            residual_xyz_td, residual_xyz_fd = total_xyz_td.copy(), total_xyz_fd
        # Only here may the residual change: one waveform/TDI sum for all accepted blocks.
        if proposals:
            parameters = np.concatenate(proposals)
            print(f'  Subtract {len(parameters)} fitted DWDs from {np.count_nonzero(accepted_counts)} blocks together.', flush=True)
            links = DWD_simulation_gpu(parameters, detector_config)
            response_model = lisagwresponse.GalacticBinary(A=1., f=.003, ra=0., dec=0.,
                orbits=str(orbit_file), orbit_interp_order=5, shift_sun='never')
            links_file = work / 'recovered_links.h5'
            with h5py.File(links_file, 'w') as saved:
                response_model.init_response_file(saved, detector_config['dt'], len(links), -2 * N_day / fs, 'tcb')
                saved['tcb/y'][:] = links
            del links
            instrument = lisainstrument.Instrument(
                **{name: value for name, value in detector_config.items()
                   if name not in ('T_obs', 'gpu_batch_size', 'minimum_envelope_samples', 'include_padding')},
                gws=str(links_file), size=N + 2 * N_day, t0=-N_day / fs)
            instrument.disable_all_noises()
            instrument.disable_dopplers()
            data = pytdi.Data.from_instrument(instrument.export_numpy_full())
            new_xyz_td = tdi_gpu(data.measurements, data.args, keep=slice(N_day, N_day + N)) / instrument.central_freq
            accumulated_xyz_td += new_xyz_td
            residual_xyz_td = total_xyz_td - accumulated_xyz_td
            residual_xyz_fd = np.fft.rfft(residual_xyz_td, axis=-1)
            del data, instrument, response_model, new_xyz_td
            accepted_parameters.append(parameters)
            accepted_doppler.extend(proposal_doppler)
            accepted_coefficients.extend(proposal_coefficients)
            accepted_ids.extend(proposal_ids)
            accepted_blocks.extend(proposal_blocks)
            last_accepted_round = round_number
            status[status == b'proposed'] = 'accepted'
\end{lstlisting}

Accepted-fit collections, accumulation and residual arrays are initialized
here on the first round, when the block pass reaches subtraction.
If no sources are accepted, the residual remains the original total observation
and the recovered signal remains zero.

The residual changes only after every selected block has been assessed.
Accepted physical parameter rows from all those blocks are concatenated and
passed to \code{DWD_simulation_gpu} once. Its summed six-link response is wrapped
in the native LISA response-file format and converted by a noise-free native
instrument. Full measurement dictionaries are then passed to \code{tdi_gpu}.

The instrument filter and moving-arm delays follow the simulation input's
configuration. One padding day is removed from each end of the instrument/TDI
stream. Division by the central optical frequency produces the fractional-
frequency XYZ convention of the stored observation.

Let \(D(t)\) denote the original three-channel observation, \(H_{r-1}(t)\) the
previously accepted waveform and \(\Delta H_r(t)\) this round's new accepted sum.
The implemented update is
\begin{align}
H_r(t)&=H_{r-1}(t)+\Delta H_r(t),\\
R_r(t)&=D(t)-H_r(t),\\
d_r&=\operatorname{rfft}(R_r).
\end{align}
Previously accepted source parameters stay frozen. The next round uses the
fresh FFT of the new shared residual. Models retain their physical sidebands
even when their carrier frequencies were selected in a restricted interval.

\Needspace{8\baselineskip}
\subsection{Step 10: record the round and decide whether to continue (lines 329--346)}


% SOURCE eric_subtraction_all.py 329 346
\begin{lstlisting}[language=Python,firstnumber=329,numbers=left,showlines=true]
        # Step 10. record the round and decide whether to continue.
        recovered_ids = np.concatenate(accepted_ids) if accepted_ids else np.empty(0, dtype=np.int64)
        round_summary = dict(block_status=status, selected_count=selected_counts, accepted_count=accepted_counts,
            recovered_dwd_id=recovered_ids)
        with h5py.File(output_file, 'a') as saved:
            group = saved[f'rounds/{label}']
            _save_tree(group, round_summary)
            group.attrs.update(complete=True, committed=bool(proposals), visited_block_count=selected_block_count,
                selected_block_count=selected_block_count,
                accepted_block_count=int(np.count_nonzero(accepted_counts)),
                residual_output_sha256=hashlib.sha256(residual_xyz_fd.tobytes()).hexdigest())
            for block_index in np.flatnonzero(accepted_counts):
                saved[f'rounds/{label}/blocks/{block_index:05d}'].attrs.update(accepted_for_subtraction=True,
                    subtraction_backend='DWD_simulation_gpu + native Instrument + tdi_gpu')
            saved.attrs.update(last_completed_round=round_number, last_accepted_round=last_accepted_round)
        if not proposals:
            stopping_reason = 'no_new_candidates' if not np.any(selected_counts) else 'no_accepted_blocks'
            break
\end{lstlisting}

The recovered-ID array is assembled here from accepted fit records, when
recording the round result.


Round statuses are stored over the whole indexed grid. Selected blocks can
lack F-statistic candidates, have a rejected fit, or commit accepted components.
Other indices remain outside the subtraction range.

If no block proposes an accepted waveform, the controller stops with
\code{no_new_candidates} or \code{no_accepted_blocks}. There is no maximum
round count, catalogue-exhaustion or ghost-fraction stopping rule.

\Needspace{8\baselineskip}
\subsection{Step 11: save the final results and finish (lines 347--365)}


% SOURCE eric_subtraction_all.py 347 365
\begin{lstlisting}[language=Python,firstnumber=347,numbers=left,showlines=true]
# Step 11. save the final results.
parameters = np.concatenate(accepted_parameters) if accepted_parameters else np.empty((0, 8))
ids = np.concatenate(accepted_ids) if accepted_ids else np.empty(0, dtype=np.int64)
final_data = dict(signal_xyz_td=accumulated_xyz_td, residual_xyz_td=residual_xyz_td,
    signal_xyz_fd=np.fft.rfft(accumulated_xyz_td, axis=-1), residual_xyz_fd=residual_xyz_fd,
    parameters=parameters, dwd_id=ids,
    doppler=np.concatenate(accepted_doppler) if accepted_doppler else np.empty((0, 4)),
    coefficients=np.concatenate(accepted_coefficients) if accepted_coefficients else np.empty((0, 4)),
    block_index=np.concatenate(accepted_blocks) if accepted_blocks else np.empty(0, dtype=np.int64))
with h5py.File(output_file, 'a') as saved:
    _save_tree(saved.create_group('final_subtraction'), final_data)
    saved['final_subtraction'].attrs['accepted_round'] = last_accepted_round
    saved.attrs.update(complete=True, single_file_complete=True, stopping_reason=stopping_reason,
        stopping_condition_met=True, iteration_stable=False)
shutil.rmtree(work)
print(f'Saved {output_file}: {round_number} global rounds, {len(parameters)} accepted components; {stopping_reason}.', flush=True)
pipeline_result = dict(output_file=str(output_file), rounds=round_number, source_count=len(parameters),
            selected_block_indices=selected_block_indices.tolist(), total_block_count=block_count,
            last_accepted_round=last_accepted_round, stopping_reason=stopping_reason)
\end{lstlisting}

The final result contains the accumulated TD and FD signal, shared residual,
accepted physical parameters, four-amplitude coefficients, Doppler points,
source IDs and global block indices. The zero-source case writes empty parameter
arrays and retains the original observation exactly.

Set \code{subtraction_frequency_range} in Step 2 to \code{(f_min, f_max)}
in Hz. The final \code{pipeline_result} dictionary records
\code{selected_block_indices}. Root completion is written after final arrays;
successful execution removes the workspace. The implementation does not
provide transactional replacement of a previous result or an automatic resume
from a partially written result.

\Needspace{8\baselineskip}
\section{Mathematical model and normalization}
\label{sec:math}

\Needspace{8\baselineskip}
\subsection{Raw FFT and the correlated XYZ covariance}

For a real channel \(x_c[m]\), the forward coefficients are
\[
d_{c,k}=\sum_{m=0}^{N-1}x_c[m]\exp(-2\pi i km/N).
\]
There is no \(1/N\) or sample-spacing factor in this forward FFT. For interior
positive-frequency bins, the one-sided CSD convention is
\begin{equation}
\mathbb E[n_k n_k^\dagger]=\frac{T_{\rm obs}f_s^2}{2}C_k,
\qquad
P_k=\frac{2}{T_{\rm obs}f_s^2}d_kd_k^\dagger.
\label{eq:rawcov}
\end{equation}
The periodogram of a single bin is rank one. Neighboring-frequency smoothing
produces the full covariance used for the likelihood.

The inspected inner product retains complex XYZ cross terms:
\begin{equation}
\ip{a}{b}_C=\frac{4}{T_{\rm obs}f_s^2}
       \operatorname{Re}\sum_k a_k^\dagger C_k^{-1}b_k.
\label{eq:inner}
\end{equation}
If \(C_k=L_kL_k^\dagger\), Cholesky whitening of each channel vector followed by
concatenation of real and imaginary parts gives an ordinary real dot product
equal to Equation~\eqref{eq:inner}. This is the convention used by the GPU
amplitude integral.

DC is excluded from the covariance estimate and retained fitting bins.
Nyquist is included in the full default grid, while the covariance and
inner-product formulas do not introduce a special real-endpoint weight.
Interior-bin normalization is the verified reference; endpoint handling remains
a specific check for searches near Nyquist, as discussed in Section~\ref{sec:findings}.

\Needspace{8\baselineskip}
\subsection{Measured covariance smoothing and peak downweighting}

The smoothing helper uses a normalized cosine kernel with half-width 256 bins,
giving 513 support entries:
\[
K_\ell\propto\tfrac12[1+\cos(\pi\ell/256)],\qquad
|\ell|\leq256,\qquad \sum_\ell K_\ell=1.
\]
This smooths the covariance in frequency, not the observation in time.
With peak downweighting, it iterates at most 12 times. For
\(s_k=\operatorname{Re}\operatorname{tr}(\widehat C_k^{-1}P_k)\), the weight is
\(w_k=\max(1-s_k/15,0)^2\). The weighted local mean is divided by its kernel
weight sum and the factor \(c=0.8503365527522067\).

For ideal proper-complex Gaussian XYZ noise, \(S\sim\operatorname{Gamma}(3,1)\).
Writing \(w(S)=(1-S/15)^2\mathbf1_{S<15}\), the scale correction is
\[
c=\frac{\mathbb E[S\,w(S)]}{3\,\mathbb E[w(S)]}.
\]
Its ideal-model derivation does not establish unbiasedness for coherent
foreground, finite windows or covariance estimated from the same samples.
The helper symmetrizes its estimate and adds a small trace-scaled diagonal
term. The full-band screen adds a further positive scale floor, or an
\(I_3\,10^{-100}\) fallback for an exactly zero chunk. Per-fit covariance is
stored, but numerical health and estimator convergence are not independently
certified by a completion flag.

\Needspace{8\baselineskip}
\subsection{Four linear amplitudes and the F projection}

At a fixed Doppler point \(\theta=(f,\dot f,\lambda,\beta)\), one source is
\[
h(\theta,a)=\sum_{\mu=1}^4 a_\mu H_\mu(\theta).
\]
The physical convention is \(A_+=A(1+\cos^2\iota)\) and
\(A_\times=2A\cos\iota\). With \(c_\psi=\cos2\psi\),
\(s_\psi=\sin2\psi\), \(c_\phi=\cos\phi_0\), \(s_\phi=\sin\phi_0\),
\begin{align}
\mathcal A_1&=A_+c_\psi c_\phi+A_\times s_\psi s_\phi,\\
\mathcal A_2&=A_+s_\psi c_\phi-A_\times c_\psi s_\phi,\\
\mathcal A_3&=-A_+c_\psi s_\phi+A_\times s_\psi c_\phi,\\
\mathcal A_4&=-A_+s_\psi s_\phi-A_\times c_\psi c_\phi.
\end{align}
The coefficient arrays use \(a_\mu=\mathcal A_\mu/10^{-22}\), with the
corresponding \(10^{-22}\) basis scale.

Define \(x_\mu=\ip{H_\mu}{d}_C\) and
\(M_{\mu\nu}=\ip{H_\mu}{H_\nu}_C\). The amplitude log likelihood ratio is
\(a^\mathsf T x-a^\mathsf T M a/2\). Its maximized value gives
\[
2\mathcal F=x^\mathsf T M^+x,
\]
where the current vectorized helper uses a Hermitian pseudoinverse with
\code{rcond=1e-12}. This admits rank-deficient bases as a projection. Numerical
agreement with the existing full-rank projection is checked separately from
selection of the recorded bank templates.

The SNR used for final acceptance is the fitted waveform norm
\[
\rho_i=\sqrt{\max\{\ip{h_i^{\rm fitted}}{h_i^{\rm fitted}}_C,0\}}.
\]
It uses fitted conditional-mean amplitudes and the measured effective
covariance. It is distinct from a stored catalogue SNR and from the conditional
F statistic after other fitted or pending components have been removed.

\Needspace{8\baselineskip}
\subsection{Joint amplitude marginalization and the zero-source model}

For \(n\) sources, stack \(4n\) basis vectors, including their cross-source
Gram terms. The normalized amplitude prior is
\(a\sim\mathcal N(0,\sigma^2 I)\), with \(\sigma=10\) in the scaled coefficient
units. Let \(Q=M+\sigma^{-2}I\). The amplitude-integrated likelihood is
\begin{equation}
\log p(d\mid\theta,n,C)
=-\tfrac12\ip{d}{d}_C+\tfrac12x^\mathsf TQ^{-1}x
 -\tfrac12\log\det Q-4n\log\sigma.
\label{eq:marginal}
\end{equation}
The determinant and prior normalization penalize additional amplitude
dimensions. For \(n=0\), the amplitude terms vanish.

The GPU integral pads to four times the configured capacity, 196 rows at the
default capacity of 49. Inactive zero-basis blocks contribute
\(Q=\sigma^{-2}I\), whose determinant and prior terms cancel. Stored amplitude
draws use the conditional Gaussian covariance \(Q^{-1}\); the subtracted
point estimate uses its conditional mean \(Q^{-1}x\).

Noise covariance and the residual are fixed within each block fit. Their
noise-normalization constant can be omitted from within-fit ratios, but
likelihood traces from different residuals or covariances are not directly
comparable evidences and should not be pooled as one posterior.

\Needspace{8\baselineskip}
\section{RJMCMC and waveform helper contracts}
\label{sec:sampler}

\Needspace{8\baselineskip}
\subsection{Signed sampled coordinates and count prior}

The global controller explicitly sets \code{chirp_coordinate='linear'}:
\begin{align}
q_f&=f_cT_{\rm obs}\ln(f/f_c),&
q_{\dot f}&=\dot f\,T_{\rm obs}^2,\\
q_\lambda&=\lambda\bmod2\pi,&q_\beta&=\sin\beta.
\end{align}
The inverse is \(f=f_c\exp[q_f/(f_cT_{\rm obs})]\),
\(\dot f=q_{\dot f}/T_{\rm obs}^2\), and \(\beta=\arcsin q_\beta\).
The helper's general default log-chirp option does not describe this caller.
Its rectangular prior implies a \(1/f\) frequency density, uniform signed
derivative, uniform longitude and isotropic sky latitude.

For capacity \(M_{\max}\), the exchangeable active-slot prior has
\[
\log p(q,\mathrm{active})
=-n\log V-\log\binom{M_{\max}}n,
\quad 0\leq n\leq M_{\max}.
\]
Summing masks and integrating active coordinates gives equal integrated count
weight for \(n=0,\ldots,M_{\max}\). The class name \code{CataloguePrior}
refers to this variable-count parameter set, not to an external source list.
The sampled count concerns additional sources in this block and round; the
cumulative recovered count can exceed 49 across blocks and rounds.

\Needspace{8\baselineskip}
\subsection{Proposals, temperatures and the selected point estimate}

Within-model Gaussian proposals move one active source in the encoded
coordinates, with standard deviations \((0.025,0.3,0.015,0.015)\) and temperature
scaling. The derivative step is linear here, rather than a multiplicative
log-chirp step. Birth proposals mix 90\% candidate-centered boxes with 10\%
full-prior draws; when no centers are supplied, the full prior is used.
The custom RJ factor compensates exchangeable slot selection. The code and
installed Eryn behavior jointly determine balance and mixing.

The four temperatures form one ladder, not four independent cold chains.
After burn-in, the helper chooses the modal cold source count, then the
highest marginal-likelihood draw at that count. Conditional-mean amplitudes
at that Doppler point generate the fitted waveform. This is a selected
point estimate; it does not average over source count, physical parameters
or subtraction uncertainty.

Stored traces are useful diagnostics, but this execution does not assess
independent-chain agreement, effective sample size or robustness to longer
sampling. A finite stopping condition or constant count trace is insufficient
to claim convergence.

\Needspace{8\baselineskip}
\subsection{DWD parameter-table GPU interface}
\label{sec:gpu_api}

\code{DWD_simulation_gpu} in \code{eric_generate_waveform.py} accepts either
an \((n,8)\) float table or a structured physical-parameter table. Additional
ID/SNR fields are ignored for synthesis. Each supplied row contributes once,
including repeated parameter sets; no LDC lookup is needed.

It returns summed float64 dimensionless link fractional-frequency data in
native link order \((12,23,31,13,32,21)\). With padding enabled, the shape is
\((N+4N_{\rm day},6)\), covering two days at each side. A noise-free instrument
then produces measurements with \(N+2N_{\rm day}\) samples. The final TDI crop
returns the central \((3,N)\) observation. The observation producer and
subtraction controller use this same sequence.

\Needspace{8\baselineskip}
\subsection{Native-expression GPU TDI interface}
\label{sec:tdi_gpu_api}

\code{tdi_gpu} defaults to \code{combination='XYZ2'}. It accepts six science-link
measurement streams, or the full native measurement dictionary. The pipelines
use dictionaries containing science, test-mass and reference beatnotes so that
the standard instrumental-noise combinations are retained. Its input must
already include instrument filtering.

The detector configuration can provide an orbit file or precomputed native
\code{Data.args} containing moving-arm delays and delay derivatives. The
implementation uses the installed PyTDI Michelson expressions with GPU
Lagrange interpolation and preserves native orders 5/31 and float64 arithmetic.
TDI itself adds no noise, FFT or file output. The output shape is
\((\ldots,3,N_{\rm kept})\) in X/Y/Z order and in the input's units.
The current-controller native comparison is recorded in
Section~\ref{sec:tdi_api_validation}.

\Needspace{8\baselineskip}
\section{Array contracts and HDF5 interpretation}
\label{sec:schema}

\Needspace{8\baselineskip}
\subsection{Array shapes and time-domain units}

{\small
\begin{longtable}{P{0.33\textwidth}P{0.25\textwidth}P{0.34\textwidth}}
\toprule Array & Shape & Scope\\\midrule\endfirsthead
\toprule Array & Shape & Scope\\\midrule\endhead
Original, accumulated signal, residual XYZ TD & \((3,N)\) & Whole observation.\\
Shared XYZ FD & \((3,\lfloor N/2\rfloor+1)\) & Raw real FFT.\\
Block measured XYZ FD & \((3,K)\) & Retained bins including sidebands.\\
Fixed block covariance & \((K,3,3)\) & Full complex XYZ CSD.\\
Source basis set & \((n,4,3,K)\) & Four amplitude bases per intrinsic point.\\
Fitted coefficients & \((n,4)\) & Amplitudes scaled by \(10^{-22}\).\\
Fitted Doppler rows & \((n,4)\) & \(f,\dot f,\lambda,\beta\).\\
Recovered physical rows & \((n,8)\) & Physical waveform parameters.\\
Summed padded GW links & \((N+4N_{\rm day},6)\) & Before instrument filtering.\\
Padded native measurements & \(N+2N_{\rm day}\) per beatnote & Before central TDI crop.\\
Source IDs, global block IDs & \((n,)\) & Separate identity fields.\\
\bottomrule
\end{longtable}}

At 360 days and \(0.1\Hz\), \(N=3{,}110{,}400\) and
\(N_{\rm day}=8{,}640\). Thus the summed padded links have 3,144,960 rows,
native measurements have 3,127,680 samples, and final XYZ has 3,110,400
samples per channel. The physical input amplitude is strain; the returned
links and normalized XYZ observation are dimensionless fractional-frequency
responses. Matching dimensions and units are both required before addition
or subtraction.

\Needspace{8\baselineskip}
\subsection{Root groups and per-round records}

\begin{description}[leftmargin=0pt,labelwidth=0pt,itemsep=5pt]
\item[\code{config/}] Frequency-block pairs and indices, subtraction range, selected
indices, thresholds, derivative bounds, sampling effort and resource limits.
\item[\code{blocks/<index>/}] Static selected-block windows, retained bins,
physical frequencies and the original full block interval.
Outside blocks need not have a group; their indices remain in
\code{config/} and the round status arrays.
\item[\code{rounds/<round>/}] Whole-grid statuses and source counts, residual input/output digests and commit state.
\item[\code{rounds/<round>/blocks/<index>/}] Fixed covariance, search result,
fitted posterior, recovery rows, post-fit labels, acceptance masks and
embedded sampler checkpoints.
\item[\code{final_subtraction/}] Full TD/FD signal and residual, cumulative
physical rows, Doppler points, amplitudes, source IDs and global block IDs.
\end{description}

Root \code{complete=True} means final publication finished.
Round \code{complete=True} means that attempt was recorded;
\code{committed=True} means it updated the residual.
\code{stopping_condition_met} excludes resource-capped termination.
\code{iteration_stable=False} records publication without a stability claim.
The per-block sampler record separately states whether independent-chain
convergence was assessed; this information is not a root attribute.

\Needspace{8\baselineskip}
\Needspace{13\baselineskip}
\subsection{A direct result reader}

\begin{lstlisting}[numbers=left,firstnumber=1]
with h5py.File('data/dwd_subtraction_all.h5', 'r') as saved:
    complete = saved.attrs['complete']
    grid_indices = saved['config/block_indices'][:]
    selected_indices = saved['config/selected_block_indices'][:]
    requested_range = saved['config/subtraction_frequency_range'][:]
    parameters = saved['final_subtraction/parameters'][:]
    source_ids = saved['final_subtraction/dwd_id'][:]
    source_blocks = saved['final_subtraction/block_index'][:]
    residual = saved['final_subtraction/residual_xyz_td'][:]
\end{lstlisting}

Calculations on these arrays belong outside the HDF5 loading scope.
A catalogue label is post-fit metadata; a negative label is not a rejection.
The original observation remains a separate input, while checkpoint arrays
needed to inspect completed fits are embedded in the result.

\Needspace{8\baselineskip}
\section{Verification and measured scope}
\label{sec:verification}
\label{sec:tdi_api_validation}

\Needspace{8\baselineskip}
\subsection{Current indexed-range controller}

Runtime validation guards and explicit exceptions have been removed from the
controller and its reviewed helpers. The numerical evidence below records the
earlier implementation; valid-input regression checks accompany this revision.

The current primary-source SHA-256 is
\path{6e34519a839f511d13ac10e370375f5d891bd4bc3710f424a20eaafd143c9db6}.
The current revision passed 43 server regressions: 19 sequential-controller
checks, five F-statistic/covariance checks, 17 observation-file checks, and two
TDI/subtraction checks. These include strict selection on saved relative SNR,
preservation of saved IDs and all 11 numerical table columns, reuse of saved
detector settings, and the native Kaiser-filter coefficient/group-delay round trip.
The numerical 11-column, named-table and four-coordinate inputs produce identical F-statistics,
amplitude coefficients and encoded sampler coordinates in the regression
fixtures; the input tables remain unchanged.

The existing SNR result was migrated to a numerical \code{catalogue} dataset
with shape \((29{,}000{,}000,11)\). A streaming comparison of every saved value
against the original catalogue rows gave identical SHA-256 digests, preserving
the parameters, both SNRs and original merged-row IDs. The revised
\code{eric_snr.py} writes this table during its final relative-SNR pass; a
fixture checks its actual export statements and saved column names. The current
subtraction reader loads 8,336 qualifying rows without reconstructing IDs.

A ten-source, one-day moving-orbit GPU smoke compared the new 11-column table
input with the existing eight-parameter input to \code{DWD_simulation_gpu}.
Both returned shape \((8{,}640,6)\) at \(0.1\Hz\), with maximum absolute
difference and relative \(L^2\) difference zero. The input table was unchanged.
This verifies the table-input route; it is separate from the native LISA
comparisons below. The records are
\path{/home/genshin/dwd_injection/tmp/review_sync_20261012.gUWiIy/merged_snr_table_verification.json},
\path{/home/genshin/dwd_injection/tmp/review_sync_20261012.gUWiIy/merged_table_integration.json}, and
\path{/home/genshin/dwd_injection/tmp/review_sync_20261012.gUWiIy/merged_table_gpu.json}.

The updated simulation was rerun on GPU1's A100 for all 29,000,000 DWDs over
360 days at \(0.1\Hz\), completing in about 19.5 minutes. The regenerated
observation stores \code{duration}, \code{rate}, and \code{detector_config}.
Its total, signal and noise TDI arrays each have shape \((3,3{,}110{,}400)\)
and contain finite values. The relative \(L^2\) difference between total and
signal plus noise is \(1.93\times10^{-14}\). The companion orbit agrees with
the embedded orbit. The run record is
\path{/home/genshin/dwd_injection/tmp/detector_config_rerun_20261012.55wr6r_k/verification.json}.

The following remote smoke and timing results describe the earlier blind-search
implementation and do not establish recovery for the current template bank.

The A100 integration smoke reused a 30-day \(0.1\Hz\) native LISA observation
with one unlisted DWD and continuous ACC+OMS noise, moving unequal arms,
and an empty SNR catalogue. Its 32 burn-in plus 64 production cycles test
integration rather than convergence.

{\small
\begin{tabularx}{\textwidth}{P{0.29\textwidth}P{0.19\textwidth}X}
\toprule Requested range (Hz) & Global indices & Observed result\\\midrule
0.002--0.00203 & 199, 200, 201, 202 & No sampler or TDI call; original observation preserved exactly.\\
0.0030005--0.003008 & 300 & One accepted source, ID \(-1\), one sampler and one TDI call; only this block visited in both global rounds.\\
\bottomrule
\end{tabularx}}

The fitted carrier was \(0.0030031746626394866\Hz\), within the requested
interval. The recovered signal's relative \(L^2\) difference from the native
injected signal was 0.0182088121. This is a fitted-waveform error, not a
GPU interpolation error. GPU TDI on the recovered measurements was bitwise
identical to native PyTDI, with maximum absolute difference zero.
Reopened files passed exact residual and FFT identities; temporary workspaces
had been removed.

\Needspace{8\baselineskip}
\subsection{Residual-only screen validation before range selection}

The earlier residual-only controller snapshot tested all 5,000 blocks on
the same 30-day fixture. Native noise alone fell below the measured spectral
threshold everywhere, requiring no sampler or TDI and taking 2.287 seconds.
One unlisted source plus native noise required two short sampler calls,
recovered one source, and took 92.790 seconds including a native TDI comparison.
The recorded earlier controller hash is
\path{19733a8fa2a424101642e4cc3182f63851e4e6e6f55efbc12dee4944ca0b4655}.
It precedes the indexed-range addition and is not the current source identity.
The evidence is \code{tmp/subtraction_blind_20261012/scientific.json}.

\Needspace{8\baselineskip}
\subsection{Forced full-pipeline noise-block benchmark}

A separate benchmark bypassed the spectral and no-seed sampler gates for
one native-noise block near \(0.003\Hz\). It searched every half-bin carrier,
32 sky positions and nine signed chirps, then ran 500 burn-in plus 1,500
production cycles with capacity 49 and four temperatures on an A100.

{\small
\begin{tabularx}{\textwidth}{P{0.30\textwidth}P{0.18\textwidth}X}
\toprule Stage & Time & Measured work\\\midrule
F-statistic search & 31.690 s & 14,976 trial templates; maximum \(F=11.083743\), no surviving seeds.\\
Unconditional RJMCMC & 70.158 s & 8,002 source-waveform evaluations and 8,004 likelihood calls.\\
Complete benchmark & 102.918 s & Includes setup, covariance, geometry and result storage; excludes Python imports and fixture generation.\\
\bottomrule
\end{tabularx}}

All 1,500 cold production draws contained zero sources, while a hotter replica
briefly visited one source. The empirical modal model therefore had nothing
to subtract. This one-block result is not a calibrated false-positive rate,
an independent-chain convergence check or an annual runtime measurement.
The evidence is \code{tmp/subtraction_empty_full_20261012/result.json}.

\Needspace{8\baselineskip}
\subsection{Scope of the native comparison}

The current-controller TDI check applies to the measurement arrays actually
generated in the bounded recovery smoke. It verifies agreement with native
PyTDI for that case. It does not establish exhaustive parity over every
possible waveform, orbit, interpolation boundary, signal population or GPU.
Earlier full-catalogue injection timings describe the observation producer;
they are not subtraction or inference throughput.

\Needspace{8\baselineskip}
\section{Review findings and scientific limits}
\label{sec:findings}

\Needspace{8\baselineskip}
\subsection{Search completeness and threshold calibration}

The finite template bank can omit an unlisted source or miss a weak, broad or crowded
source. The existing high-SNR single-source smoke does not calibrate recovery
near fitted SNR 7. Template maximization,
same-data covariance estimate and repeated global rounds jointly determine
false alarms and misses. Validate controlled noise ensembles and injected
populations across frequency, chirp, sky, crowding and range boundaries.

\Needspace{8\baselineskip}
\subsection{Endpoint and covariance health}

Full-band execution reaches Nyquist, whereas the generic positive-frequency
normalization and spectral reference use interior-bin assumptions. A dedicated
real-endpoint comparison is needed before assigning the same statistical
interpretation there. Save or independently inspect covariance eigenvalues,
conditioning and estimator stability when evaluating difficult bands.

\Needspace{8\baselineskip}
\subsection{Shared rounds remain conditional fits}

Each block has its own conditional posterior against a common snapshot.
Previously accepted sources are frozen. Later acceptance removes already
pending components, so proposal order can affect a crowded shared feature.
The resulting global subtraction is coordinated across the selected blocks,
but it is not one simultaneous posterior over the entire accepted catalogue.
Overlapping support regions and multiple close sources require dedicated
recovery and uncertainty tests.

\Needspace{8\baselineskip}
\subsection{FD proposal assessment and native TD subtraction}

Fitting and acceptance use an FD response model. Accepted physical rows are
then synthesized, instrument-filtered and passed through moving-arm TDI in
the time domain. Numerical agreement between these routes must be measured
over the intended duration, frequency and chirp range. The current smoke
checks the final residual identity and a bounded recovery error; it does not
bound every FD-to-TD modelling discrepancy or guarantee a global likelihood
increase in a full-year crowded fit.

\Needspace{8\baselineskip}
\subsection{Reproducibility, labels and failure recovery}

The output stores covariance, configuration, sample arrays and checkpoints.
The sampler's separate NumPy amplitude generator and the complete software
environment are not fully reconstructible from the controller configuration.
Post-fit catalogue association is greedy and can alter labels, but it does
not change the data-based acceptance predicate.

Opening the final path with \code{'w'} and incremental round writes leaves
a partial result if execution fails; old output is already replaced.
Temporary files remain available for diagnosis on an exception, but automatic
resume or transactional publication is not implemented. Preserve successful
results deliberately when rerunning a new experiment.

\Needspace{8\baselineskip}
\subsection{Stopping and residual-background interpretation}

\code{no_new_candidates} and \code{no_accepted_blocks} refer to the selected
range under the configured search and acceptance rules. They do not prove
all resolvable DWDs have been removed, nor infer posterior convergence.
A residual used for foreground or stochastic-background analysis also contains
missed signals, waveform mismatch, noise-estimation effects and subtraction
uncertainty. A single conditional-mean subtraction does not propagate the
posterior uncertainty of the recovered source population.

\Needspace{8\baselineskip}
\section{Performance and practical use}
\label{sec:performance}

The full FFT and covariance estimate run once per global round. Every selected
block evaluates its recorded templates. RJMCMC and posterior storage follow
when templates pass the F-statistic threshold. Response geometry is
shared. Accepted sources from all selected blocks are synthesized and passed
through native instrument filtering and GPU TDI once per committed round.

GPU synthesis and whitened response/integral evaluation accelerate numerical
work. Catalogue scans, covariance smoothing, proposal/state bookkeeping,
conditional amplitude draws and HDF5 I/O still involve the CPU. The pipeline
is GPU assisted rather than entirely GPU resident.

At 360 days and \(0.1\Hz\), one float64 XYZ array is about 71.2 MiB, one
complex128 real-FFT XYZ array about 71.2 MiB, and one six-link padded array
about 144.0 MiB. Concurrent arrays, geometry, waveform batches, covariance
and sampler state add to these base sizes. After the first attempted fit,
sorting 29 million catalogue frequencies and row indices also requires
hundreds of MiB of host memory, even though that scan supplies labels only.

These storage scales are deterministic array estimates. The measured
30-day benchmark and short integration smokes cannot be multiplied into a
validated annual production-runtime claim.

\Needspace{8\baselineskip}
\section{Scope, source identity and coverage}
\label{sec:scope}

The primary local file is
\path{/Users/liangqingyuan/Desktop/codex/DWD injection/eric_subtraction_all.py}.
Its server counterpart is
\path{/home/genshin/dwd_injection/eric_subtraction_all.py}.
The walkthrough reproduces every line of the 365-line sequential Python script
in its actual source order. Imports and parameter assignments are part of the
Python file itself. Code text, indentation and line numbers match the source
exactly; no separate presentation numbering or generated parameter block is
used. Steps 4--10 repeat for each global round. Appendix~\ref{app:source}
reproduces the full file, and Appendix~\ref{app:coverage} maps its step ranges.

The review is tied to the following source identities. Source hashes identify
the inspected implementation; the bounded evidence in
Section~\ref{sec:verification} defines the tested numerical scope.


\textbf{\code{eric_subtraction_all.py}} (365 lines)\\\path{6e34519a839f511d13ac10e370375f5d891bd4bc3710f424a20eaafd143c9db6}


\textbf{\code{eric_simulation_DWD.py}} (134 lines)\\\path{f902efc3b63c17772d59472e1c912c134ff9f38456304d8b1157a920ab6159f9}


\textbf{\code{eric_snr.py}} (360 lines)\\\path{7441c8ba7b7e0dc4680e5e4a14b6bb435be07684c68816f300911d46f1b0510e}


\textbf{\code{eric_generate_waveform.py}} (266 lines)\\\path{ab9c96bbfa3de392cac9d34034aac56ed5cda45770f544f655d13d5df0a4eceb}


\textbf{\code{eric_rjmcmc.py}} (316 lines)\\\path{9401a9332332010313416a384b44c89ca49bea7f95860e37247d8f9a3e934a46}


\textbf{\code{script/dwd_blind_search.py}} (68 lines)\\\path{5c4cdbb56c662195ba40b20e3802c40a1e6ca9a9c66075de8618ac1f7fe3c7b0}


\appendix
\Needspace{8\baselineskip}
\section{Complete primary-source listing}
\label{app:source}

This is the complete current \mainfile, with its original line numbers.


% SOURCE eric_subtraction_all.py 1 365
\begin{lstlisting}[language=Python,firstnumber=1,basicstyle=\ttfamily\scriptsize,numbers=left,showlines=true]
'''Step 1: import all the necessary packages.'''
from pathlib import Path
from tempfile import mkdtemp
import hashlib
import json
import shutil
import h5py
import numpy as np
import lisagwresponse
import lisainstrument
import pytdi
from eric_generate_waveform import DWD_simulation_gpu, tdi_gpu
from script.eric_Fstats import (
    build_response, physical_amplitude_coordinates,
    physical_amplitudes, reconstruct_waveform, noise_inner_product,
)
from script.dwd_blind_search import (
    residual_covariance, template_f_statistic_search, amplitude_f_statistic,
)
from script.dwd_bank_association import assign_bank_ids
from eric_rjmcmc import run_eryn_rjmcmc
from script.oneblock_observation import read_checkpoint, write_checkpoint
from script.tool import _save_tree, single_cpu

'''Step 2: define the parameters.'''
orbit_file = Path('./data/dwd_observation_full_catalogue.orbits.h5')
snr_file = Path('./data/ldc_dwd_data.h5')
observation_file = Path('./data/dwd_observation_full_catalogue.h5')
output_file = Path('./data/dwd_subtraction_all.h5')
frequency_range = (0.0001, 0.05)
f_derivative_pn = [-25, -10]
freq_block_width = 1e-05
subtraction_frequency_range = None
frequency_derivative_bounds = 10. ** f_derivative_pn[1] * np.array([-1., 1.])
F_threshold = 16.0
SNR_r_threshold = 7.0
sample_number = (500, 1500)
maximum_sources = 64
with h5py.File(observation_file, 'r') as saved:
    T_obs = saved['observation'].attrs['duration']
    fs = saved['observation'].attrs['rate']
    total_xyz_td = saved['observation/total_xyz_td'][:]
    detector_config = json.loads(saved['observation'].attrs['detector_config'])
detector_config['orbits'] = str(orbit_file)
N, N_day = int(round(T_obs * fs)), int(round(86400 * fs))
frequency_blocks = np.arange(
    frequency_range[0], frequency_range[1] - freq_block_width * 1.e-10, freq_block_width,
)[:, None] + np.array([0., freq_block_width])
frequency_blocks[-1, 1] = frequency_range[1]
block_count = len(frequency_blocks)
block_indices = np.arange(block_count, dtype=np.int64)
subtraction_range = np.array([frequency_range[0], frequency_range[1]]) if subtraction_frequency_range is None else np.asarray(subtraction_frequency_range, dtype=float)
selected_block_indices = block_indices[
    (frequency_blocks[:, 0] >= subtraction_range[0] - freq_block_width * 1.e-10)
    & (frequency_blocks[:, 1] <= subtraction_range[1] + freq_block_width * 1.e-10)]
selected_block_count = len(selected_block_indices)
fields = dict(A=1, f=2, df=3, phi0=4, iota=5, longitude=6, latitude=7, psi_bme=8)
# Step 3. read the candidate bank and prepare shared state.
with h5py.File(snr_file, 'r') as saved:
    template_dwd = saved['catalogue'][saved['catalogue'][:, 10] > SNR_r_threshold]
total_xyz_fd, fft_frequency = np.fft.rfft(total_xyz_td, axis=-1), np.fft.rfftfreq(N, d=1 / fs)
work = Path(mkdtemp(prefix='.dwd_all_', dir=output_file.parent)).resolve()
configuration = dict(duration=T_obs, rate=fs, frequency_range=np.array([frequency_range[0], frequency_range[1]]),
    frequency_blocks_hz=frequency_blocks, block_indices=block_indices, frequency_block_width=freq_block_width,
    subtraction_frequency_range=subtraction_range, selected_block_indices=selected_block_indices,
    derivative_bounds=frequency_derivative_bounds,
    chirp_coordinate='linear', f_threshold=F_threshold,
    relative_snr_threshold=SNR_r_threshold, sample_number=np.asarray(sample_number),
    minimum_sampler_capacity=maximum_sources)
with h5py.File(output_file, 'w') as saved:
    _save_tree(saved.create_group('config'), configuration)
    saved.create_group('blocks')
    saved.create_group('rounds')
# Step 4. start a global round and estimate covariance.
round_number, last_accepted_round = 0, 0
print(f'Subtraction range {subtraction_range[0]:.8g}--{subtraction_range[1]:.8g} Hz: '
      f'selected {selected_block_count} of {block_count} blocks, indices '
      f'{selected_block_indices[0]}--{selected_block_indices[-1]} (zero-based).', flush=True)
with single_cpu():
    while True:
        round_number += 1
        label = f'{round_number:02d}'
        print(f'Global round {round_number}: visit {selected_block_count} selected frequency blocks.', flush=True)
        # This complete Fourier snapshot is unchanged throughout the block loop.
        round_residual_fd = total_xyz_fd if round_number == 1 else residual_xyz_fd
        round_residual_fd.setflags(write=False)
        round_digest = hashlib.sha256(round_residual_fd.tobytes()).hexdigest()
        covariance_full = residual_covariance(round_residual_fd, T_obs, fs)
        block_order = selected_block_indices
        status = np.full(block_count, 'outside_subtraction_range', dtype='S32')
        status[selected_block_indices] = 'pending'
        with h5py.File(output_file, 'a') as saved:
            group = saved['rounds'].create_group(label)
            group.create_group('blocks')
            group.attrs.update(complete=False, committed=False, residual_input_sha256=round_digest)
        # Step 5. evaluate each selected block against its recorded templates.
        selected_counts = np.zeros(block_count, dtype=np.int64)
        accepted_counts = np.zeros(block_count, dtype=np.int64)
        proposals, proposal_doppler, proposal_coefficients, proposal_ids, proposal_blocks = [], [], [], [], []
        for block_index in block_order:
            key = f'{block_index:05d}'
            block_range = frequency_blocks[block_index]
            center = float(block_range.mean())
            block_derivative_bounds = np.asarray(frequency_derivative_bounds, dtype=float)
            padding = block_range[1] * 1.e-4 + max(abs(block_derivative_bounds)) * T_obs + 8 / T_obs
            retained_bins = np.flatnonzero((fft_frequency >= block_range[0] - padding)
                                          & (fft_frequency <= block_range[1] + padding) & (fft_frequency > 0))
            center_bin = int(round(center * T_obs))
            wide_bins = np.arange(max(1, min(center_bin - 2048, retained_bins[0] - 256)),
                min(len(fft_frequency), max(center_bin + 2049, retained_bins[-1] + 257)))
            core_indices = retained_bins - wide_bins[0]
            physical_frequency, wide_frequency = fft_frequency[retained_bins], fft_frequency[wide_bins]
            if round_number == 1 and block_index == block_order[0]:
                builder, response_times, response_geometry = build_response(work, wide_frequency, T_obs, fs, orbit_file=orbit_file)
                builder_type = type(builder)
            else:
                builder = builder_type(np.broadcast_to(np.eye(3), (len(wide_frequency), 3, 3)),
                    np.zeros((3, len(wide_frequency)), complex), response_times, response_geometry, wide_frequency, T_obs, fs)
            template_data = dict(derivative_bounds=block_derivative_bounds,
                retained_bins=retained_bins, wide_bins=wide_bins, core_indices=core_indices,
                block_index=block_index, frequency_range=block_range,
                physical_frequency_hz=physical_frequency)
            with h5py.File(output_file, 'a') as saved:
                if key not in saved['blocks']:
                    _save_tree(saved['blocks'].create_group(key), template_data)
            covariance = covariance_full[retained_bins].copy()
            inverse = np.linalg.inv(covariance)
            covariance.setflags(write=False)
            data = round_residual_fd[:, retained_bins]
            target = (template_dwd[:, fields['f']] >= block_range[0]) & (
                (template_dwd[:, fields['f']] < block_range[1]) |
                ((block_index == block_count - 1) & (template_dwd[:, fields['f']] == block_range[1])))
            search = template_f_statistic_search(data, covariance, builder, core_indices, T_obs, fs,
                template_dwd[target], f_threshold=F_threshold)
            search['template_dwd_id'] = template_dwd[target, 0].astype(np.int64)
            selected = search['template_dwd']
            selected_counts[block_index] = len(selected)
            block_data = dict(fixed_covariance=covariance,
                screening=dict(search=search, residual_input_xyz_fd=data))
            block_attributes = dict(residual_input_sha256=round_digest,
                block_index=block_index,
                fixed_covariance_sha256=hashlib.sha256(covariance.tobytes()).hexdigest(),
                proposed_for_subtraction=False, accepted_for_subtraction=False)
            status[block_index] = 'no_new_candidates'
            # Step 6. run RJMCMC and assess fitted sources.
            if len(selected):
                print(f'  Block {key}, {block_range[0]:.8f}--{block_range[1]:.8f} Hz: fit {len(selected)} candidates.', flush=True)
                checkpoints = work / f'round_{label}_block_{key}'
                posterior = run_eryn_rjmcmc(data, covariance, selected, block_range,
                    block_derivative_bounds, center, T_obs, fs, checkpoints,
                    chain_count=1, inverse_temperatures=(1., .8, .6, .4),
                    iterations=sum(sample_number), burn_in=sample_number[0],
                    maximum_sources=max(maximum_sources, len(selected)), chirp_coordinate='linear',
                    response_times=response_times, response_geometry=response_geometry, physical_frequency=physical_frequency)
                doppler, coefficients = posterior['doppler'], posterior['fit_result']['coefficients']
                new_bases = builder.build(doppler)[:, :, :, core_indices]
                new_sources = reconstruct_waveform(coefficients, new_bases)
                source_snr = np.sqrt(np.maximum(noise_inner_product(new_sources, new_sources, inverse, T_obs, fs), 0.))
                assessment_data = data.copy()
                # Fits share the round snapshot, but acceptance is conditional on
                # already proposed components to avoid subtracting a sideband twice.
                if proposal_doppler:
                    points, amplitudes = np.concatenate(proposal_doppler), np.concatenate(proposal_coefficients)
                    nearby = abs(points[:, 0] - center) <= .5 * np.ptp(block_range) + 2 * padding
                    if np.any(nearby):
                        pending = reconstruct_waveform(amplitudes[nearby],
                            builder.build(points[nearby])[:, :, :, core_indices])
                        assessment_data -= pending.sum(axis=0)
                conditional_f = np.zeros(len(doppler))
                for i in range(len(doppler)):
                    conditional_data = assessment_data - new_sources.sum(axis=0) + new_sources[i]
                    conditional_f[i] = amplitude_f_statistic(conditional_data, new_bases[i:i + 1],
                        covariance, T_obs, fs)['two_f'][0] / 2
                owned = (doppler[:, 0] >= block_range[0]) & ((doppler[:, 0] < block_range[1])
                    | ((block_range[1] == frequency_range[1]) & (doppler[:, 0] == frequency_range[1])))
                duplicate = np.zeros(len(doppler), dtype=bool)
                previous_doppler = proposal_doppler if round_number == 1 else accepted_doppler + proposal_doppler
                previous_coefficients = proposal_coefficients if round_number == 1 else accepted_coefficients + proposal_coefficients
                if previous_doppler:
                    points, amplitudes = np.concatenate(previous_doppler), np.concatenate(previous_coefficients)
                    nearby = abs(points[:, 0] - center) <= .5 * np.ptp(block_range) + 2 * padding
                    if np.any(nearby):
                        previous = reconstruct_waveform(amplitudes[nearby],
                            builder.build(points[nearby])[:, :, :, core_indices])
                        overlap = noise_inner_product(new_sources[:, None], previous[None, :], inverse, T_obs, fs)
                        norm = np.sqrt(np.maximum(noise_inner_product(new_sources, new_sources, inverse, T_obs, fs)[:, None]
                            * noise_inner_product(previous, previous, inverse, T_obs, fs)[None, :], 1.e-300))
                        duplicate |= np.any(overlap / norm > .9, axis=1)
                retained = owned & ~duplicate & np.isfinite(source_snr) & (source_snr > SNR_r_threshold) & (conditional_f >= F_threshold)
                # Deduplicate the fitted components using measured waveform overlap.
                for i in np.argsort(source_snr)[::-1]:
                    if not retained[i]:
                        continue
                    for j in np.flatnonzero(retained):
                        if j == i or source_snr[j] > source_snr[i] or (source_snr[j] == source_snr[i] and j < i):
                            continue
                        overlap = noise_inner_product(new_sources[i], new_sources[j], inverse, T_obs, fs)
                        if overlap / max(source_snr[i] * source_snr[j], 1.e-300) > .9:
                            retained[j], duplicate[j] = False, True
                retained_sum = new_sources[retained].sum(axis=0)
                improvement = float(noise_inner_product(assessment_data, assessment_data, inverse, T_obs, fs)
                    - noise_inner_product(assessment_data - retained_sum, assessment_data - retained_sum, inverse, T_obs, fs))
                accepted = np.any(retained) and improvement > 0
                status[block_index] = 'proposed' if accepted else 'fit_rejected_by_data'
                # Step 7. associate recovered sources with catalogue labels.
                # Catalogue reads happen after inference and have no acceptance role.
                if round_number == 1 and np.count_nonzero(selected_counts) == 1:
                    with h5py.File(snr_file, 'r') as saved:
                        truth_frequency = saved['catalogue'][:, fields['f']]
                        chirp = saved['catalogue'][:, fields['df']]
                    maximum_drift = float(np.max(abs(chirp) * T_obs + abs(truth_frequency) * 1.e-4 + 8 / T_obs))
                    truth_order = np.argsort(truth_frequency, kind='stable')
                    truth_frequency = truth_frequency[truth_order]
                # Truth is consulted only after fitting, and cached on disk for this block.
                truth_cache = work / f'truth_{key}.h5'
                if truth_cache.exists():
                    with h5py.File(truth_cache, 'r') as saved:
                        truth_ids = saved['dwd_id'][:]
                        truth_waveforms = saved['source_xyz_fd'][:]
                else:
                    truth_low, truth_high = block_range + np.array([-2304., 2304.]) / T_obs
                    first, last = np.searchsorted(truth_frequency, [truth_low - maximum_drift, truth_high + maximum_drift])
                    nearby_indices = np.sort(truth_order[first:last])
                    with h5py.File(snr_file, 'r') as saved:
                        records = saved['catalogue'][nearby_indices]
                    frequency = records[:, fields['f']]
                    end = frequency + records[:, fields['df']] * T_obs
                    margin = np.maximum(frequency, end) * 1.e-4 + 8 / T_obs
                    nearby = (np.minimum(frequency, end) - margin <= truth_high) & (np.maximum(frequency, end) + margin >= truth_low)
                    truth_rows = records[nearby]
                    truth_ids = truth_rows[:, 0].astype(np.int64)
                    truth_doppler = truth_rows[:, [fields[name] for name in ('f', 'df', 'longitude', 'latitude')]]
                    cos_iota = np.cos(truth_rows[:, fields['iota']])
                    truth_coefficients = physical_amplitude_coordinates(truth_rows[:, fields['A']] * (1 + cos_iota**2),
                        2 * truth_rows[:, fields['A']] * cos_iota, truth_rows[:, fields['psi_bme']], truth_rows[:, fields['phi0']]) / 1.e-22
                    truth_waveforms = np.empty((len(truth_ids), 3, len(retained_bins)), dtype=complex)
                    for first in range(0, len(truth_ids), 32):
                        truth_waveforms[first:first + 32] = reconstruct_waveform(truth_coefficients[first:first + 32],
                            builder.build(truth_doppler[first:first + 32])[:, :, :, core_indices])
                    with h5py.File(truth_cache, 'w') as saved:
                        saved.create_dataset('dwd_id', data=truth_ids)
                        saved.create_dataset('source_xyz_fd', data=truth_waveforms)
                truth_association = assign_bank_ids(new_sources, truth_waveforms, truth_ids, inverse, T_obs, fs, minimum_overlap=.9)
                matched = truth_association['matched_dwd_id'].copy()
                physical = physical_amplitudes(coefficients)
                parameters = np.column_stack((physical['A'], doppler[:, 0], doppler[:, 1], physical['phi0'],
                    physical['iota'], doppler[:, 2], doppler[:, 3], physical['psi_bme']))
                block_data.update(catalogue_association=dict(matched_dwd_id=matched,
                    best_overlap=truth_association['best_overlap'], scope='Post-fit labels only; unmatched fits may be subtracted.'),
                    # Step 8. record the accepted fits.
                    acceptance=dict(source_snr=source_snr, conditional_f=conditional_f,
                        owned_by_block=owned, duplicate_previous_or_this_round=duplicate,
                        retained=retained, likelihood_improvement=improvement),
                    recovery=dict(doppler=doppler, coefficients=coefficients, parameters=parameters, source_xyz_fd=new_sources),
                    rjmcmc={name: value for name, value in posterior.items() if name not in ('bases', 'fit_result', 'checkpoint_file')})
                block_attributes.update(proposed_for_subtraction=accepted,
                    source_xyz_fd_scope='Frequency-domain fitting and post-fit association only.')
                if accepted:
                    retained_ids = matched[retained].copy()
                    if round_number == 1 and not np.any(accepted_counts):
                        next_unknown_id = -1
                    if round_number == 1:
                        used_labels = np.concatenate(proposal_ids) if proposal_ids else np.empty(0, dtype=np.int64)
                    else:
                        used_labels = np.concatenate(accepted_ids + proposal_ids)
                    for i in range(len(retained_ids)):
                        if retained_ids[i] < 0 or retained_ids[i] in used_labels or retained_ids[i] in retained_ids[:i]:
                            retained_ids[i] = next_unknown_id
                            next_unknown_id -= 1
                    proposals.append(parameters[retained])
                    proposal_doppler.append(doppler[retained])
                    proposal_coefficients.append(coefficients[retained])
                    proposal_ids.append(retained_ids)
                    proposal_blocks.append(np.full(np.count_nonzero(retained), block_index, dtype=np.int64))
                    accepted_counts[block_index] = np.count_nonzero(retained)
                snapshots = {}
                for path in sorted(checkpoints.glob('*.h5')):
                    with h5py.File(path, 'r') as saved:
                        snapshots[path.stem] = read_checkpoint(saved)
                with h5py.File(output_file, 'a') as saved:
                    group = saved[f'rounds/{label}/blocks'].create_group(key)
                    _save_tree(group, block_data)
                    group.attrs.update(block_attributes)
                    checkpoint_group = group.create_group('checkpoints')
                    for name, snapshot in snapshots.items():
                        write_checkpoint(checkpoint_group.create_group(name), snapshot)
                    group['rjmcmc'].attrs['checkpoint_group'] = f'/rounds/{label}/blocks/{key}/checkpoints'
            else:
                with h5py.File(output_file, 'a') as saved:
                    group = saved[f'rounds/{label}/blocks'].create_group(key)
                    _save_tree(group, block_data)
                    group.attrs.update(block_attributes)
        # Step 9. subtract the combined time-domain TDI signal.
        if round_number == 1:
            accepted_parameters, accepted_doppler, accepted_coefficients, accepted_ids, accepted_blocks = [], [], [], [], []
            accumulated_xyz_td = np.zeros_like(total_xyz_td)
            residual_xyz_td, residual_xyz_fd = total_xyz_td.copy(), total_xyz_fd
        # Only here may the residual change: one waveform/TDI sum for all accepted blocks.
        if proposals:
            parameters = np.concatenate(proposals)
            print(f'  Subtract {len(parameters)} fitted DWDs from {np.count_nonzero(accepted_counts)} blocks together.', flush=True)
            links = DWD_simulation_gpu(parameters, detector_config)
            response_model = lisagwresponse.GalacticBinary(A=1., f=.003, ra=0., dec=0.,
                orbits=str(orbit_file), orbit_interp_order=5, shift_sun='never')
            links_file = work / 'recovered_links.h5'
            with h5py.File(links_file, 'w') as saved:
                response_model.init_response_file(saved, detector_config['dt'], len(links), -2 * N_day / fs, 'tcb')
                saved['tcb/y'][:] = links
            del links
            instrument = lisainstrument.Instrument(
                **{name: value for name, value in detector_config.items()
                   if name not in ('T_obs', 'gpu_batch_size', 'minimum_envelope_samples', 'include_padding')},
                gws=str(links_file), size=N + 2 * N_day, t0=-N_day / fs)
            instrument.disable_all_noises()
            instrument.disable_dopplers()
            data = pytdi.Data.from_instrument(instrument.export_numpy_full())
            new_xyz_td = tdi_gpu(data.measurements, data.args, keep=slice(N_day, N_day + N)) / instrument.central_freq
            accumulated_xyz_td += new_xyz_td
            residual_xyz_td = total_xyz_td - accumulated_xyz_td
            residual_xyz_fd = np.fft.rfft(residual_xyz_td, axis=-1)
            del data, instrument, response_model, new_xyz_td
            accepted_parameters.append(parameters)
            accepted_doppler.extend(proposal_doppler)
            accepted_coefficients.extend(proposal_coefficients)
            accepted_ids.extend(proposal_ids)
            accepted_blocks.extend(proposal_blocks)
            last_accepted_round = round_number
            status[status == b'proposed'] = 'accepted'
        # Step 10. record the round and decide whether to continue.
        recovered_ids = np.concatenate(accepted_ids) if accepted_ids else np.empty(0, dtype=np.int64)
        round_summary = dict(block_status=status, selected_count=selected_counts, accepted_count=accepted_counts,
            recovered_dwd_id=recovered_ids)
        with h5py.File(output_file, 'a') as saved:
            group = saved[f'rounds/{label}']
            _save_tree(group, round_summary)
            group.attrs.update(complete=True, committed=bool(proposals), visited_block_count=selected_block_count,
                selected_block_count=selected_block_count,
                accepted_block_count=int(np.count_nonzero(accepted_counts)),
                residual_output_sha256=hashlib.sha256(residual_xyz_fd.tobytes()).hexdigest())
            for block_index in np.flatnonzero(accepted_counts):
                saved[f'rounds/{label}/blocks/{block_index:05d}'].attrs.update(accepted_for_subtraction=True,
                    subtraction_backend='DWD_simulation_gpu + native Instrument + tdi_gpu')
            saved.attrs.update(last_completed_round=round_number, last_accepted_round=last_accepted_round)
        if not proposals:
            stopping_reason = 'no_new_candidates' if not np.any(selected_counts) else 'no_accepted_blocks'
            break
# Step 11. save the final results.
parameters = np.concatenate(accepted_parameters) if accepted_parameters else np.empty((0, 8))
ids = np.concatenate(accepted_ids) if accepted_ids else np.empty(0, dtype=np.int64)
final_data = dict(signal_xyz_td=accumulated_xyz_td, residual_xyz_td=residual_xyz_td,
    signal_xyz_fd=np.fft.rfft(accumulated_xyz_td, axis=-1), residual_xyz_fd=residual_xyz_fd,
    parameters=parameters, dwd_id=ids,
    doppler=np.concatenate(accepted_doppler) if accepted_doppler else np.empty((0, 4)),
    coefficients=np.concatenate(accepted_coefficients) if accepted_coefficients else np.empty((0, 4)),
    block_index=np.concatenate(accepted_blocks) if accepted_blocks else np.empty(0, dtype=np.int64))
with h5py.File(output_file, 'a') as saved:
    _save_tree(saved.create_group('final_subtraction'), final_data)
    saved['final_subtraction'].attrs['accepted_round'] = last_accepted_round
    saved.attrs.update(complete=True, single_file_complete=True, stopping_reason=stopping_reason,
        stopping_condition_met=True, iteration_stable=False)
shutil.rmtree(work)
print(f'Saved {output_file}: {round_number} global rounds, {len(parameters)} accepted components; {stopping_reason}.', flush=True)
pipeline_result = dict(output_file=str(output_file), rounds=round_number, source_count=len(parameters),
            selected_block_indices=selected_block_indices.tolist(), total_block_count=block_count,
            last_accepted_round=last_accepted_round, stopping_reason=stopping_reason)
\end{lstlisting}

\Needspace{30\baselineskip}
\section{Complete primary-file coverage map}
\label{app:coverage}

{\small
\begin{longtable}{P{0.065\textwidth}P{0.18\textwidth}P{0.245\textwidth}P{0.40\textwidth}}
\toprule Step & Walkthrough lines & Source lines & Explanation\\\midrule\endfirsthead
\toprule Step & Walkthrough lines & Source lines & Explanation\\\midrule\endhead
1 & 1--24 & 1--24 & Import the required packages.\\
2 & 25--57 & 25--57 & Define parameters, read the observation settings and select blocks.\\
3 & 58--73 & 58--73 & Read the recorded template bank and prepare shared state.\\
4 & 74--95 & 74--95 & Start a global round and estimate covariance.\\
5 & 96--144 & 96--144 & Evaluate each selected block against its recorded templates.\\
6 & 145--204 & 145--204 & Run RJMCMC and assess fitted sources.\\
7 & 205--249 & 205--249 & Associate catalogue labels after fitting.\\
8 & 250--292 & 250--292 & Record accepted fits and source IDs.\\
9 & 293--328 & 293--328 & Subtract the combined time-domain TDI signal.\\
10 & 329--346 & 329--346 & Record the round and decide whether to continue.\\
11 & 347--365 & 347--365 & Save the final results and remove temporary files.\\
\bottomrule
\end{longtable}}

\Needspace{8\baselineskip}
\section{Dependency source listings}
\label{app:helpers}

The following files support the reviewed execution. Their line numbers belong
to the named file and do not continue the primary controller's line numbers.

\Needspace{8\baselineskip}
\subsection{Observation producer: eric\_simulation\_DWD.py}


% SOURCE eric_simulation_DWD.py 1 134
\begin{lstlisting}[language=Python,firstnumber=1,basicstyle=\ttfamily\scriptsize,numbers=left]
# Step 1. Import packages.
import os
for thread_variable in ('OMP_NUM_THREADS', 'OPENBLAS_NUM_THREADS', 'MKL_NUM_THREADS', 'NUMEXPR_NUM_THREADS'):
    os.environ[thread_variable] = '1'
import importlib.metadata
import json
from pathlib import Path
import shutil
import tempfile
import time
import h5py
import lisaorbits
import lisagwresponse
import lisainstrument
import numpy as np
import pytdi
from script.oneblock_observation import pack_observation
from eric_generate_waveform import DWD_simulation_gpu, tdi_gpu


# Step 2. Configure parameter settings.
catalogue_file = Path('./data/LDC2_sangria_training_v2.h5')
observation_file = Path('./data/dwd_observation_full_catalogue.h5')
T_obs = 360. * 86400.
fs = 0.1
L = 2.5e9
N_acc = 3.0e-15
delta_x = 15.0e-12
gpu_batch_size = 10000
minimum_envelope_samples = 512
parameter_fields = dict(
    A='Amplitude', f='Frequency', df='FrequencyDerivative', phi0='InitialPhase',
    iota='Inclination', longitude='EclipticLongitude', latitude='EclipticLatitude',
    psi_bme='Polarization',
)


workspace = Path(tempfile.mkdtemp(prefix='.dwd_simulation_', dir=observation_file.parent))
orbit_file = workspace / 'orbit.h5'
N = int(round(T_obs * fs))
N_day = int(round(86400 * fs))
detector_config = dict(
    T_obs=T_obs, gpu_batch_size=gpu_batch_size,
    minimum_envelope_samples=minimum_envelope_samples, include_padding=True,
    dt=1 / fs, orbits=str(orbit_file), orbit_dataset='tcb/ltt',
    aafilter=('kaiser', 240., .275 * fs, .725 * fs),
    clock_offsets=0., clock_freqoffsets=0., clock_freqlindrifts=0., clock_freqquaddrifts=0.,
)

# Step 3. Simulate instrumental noise with the LISA packages.
lisaorbits.KeplerianOrbits(L=L).write(
    str(orbit_file), t0=-20e6, dt=86400., size=825, mode='w',
)
print('Simulating one continuous ACC+OMS noise realization.', flush=True)
noise_instrument = lisainstrument.Instrument(
    **{key: value for key, value in detector_config.items()
       if key not in ('T_obs', 'gpu_batch_size', 'minimum_envelope_samples', 'include_padding')}, size=N + 2 * N_day,
    t0=-N_day / fs,
    testmass_asds=N_acc, oms_asds=(delta_x, 0., 0., 0., 0., 0.),
)
noise_instrument.disable_all_noises(excluding={'test-mass', 'oms'})
noise_instrument.disable_dopplers()
noise_data = pytdi.Data.from_instrument(
    noise_instrument.export_numpy_full(),
)

# Step 4. Simulate waveforms of all DWDs on the GPU.
with h5py.File(catalogue_file) as source:
    dgb = source['sky/dgb/cat'][:]
    igb = source['sky/igb/cat'][:]

source_count = len(dgb) + len(igb)
catalogue = np.empty(source_count, dtype=[('ID', 'i8')] + [
    (field, 'f8') for field in parameter_fields.values()
])
catalogue['ID'] = np.arange(1, source_count + 1, dtype=np.int64)
for field in parameter_fields.values():
    catalogue[field] = np.concatenate((dgb[field][:, 0], igb[field][:, 0]))
del dgb, igb

with h5py.File(workspace / 'population.h5', 'x') as saved:
    saved.create_dataset('catalogues/all', data=catalogue, chunks=True)

links_file = workspace / 'population_links.h5'
print(f'Simulating all {source_count:,} DWDs with DWD_simulation_gpu.', flush=True)
started = time.perf_counter()
links = DWD_simulation_gpu(catalogue, detector_config)
waveform_runtime = time.perf_counter() - started
print(f'GPU injection completed: {source_count:,} DWDs in {waveform_runtime:.1f}s.', flush=True)
del catalogue
response_model = lisagwresponse.GalacticBinary(
    A=1., f=.003, ra=0., dec=0., orbits=str(orbit_file),
    orbit_interp_order=5, shift_sun='never',
)
with h5py.File(links_file, 'x') as saved:
    response_model.init_response_file(saved, detector_config['dt'], len(links), -2 * N_day / fs, 'tcb')
    saved['tcb/y'][:] = links
del links

# Convert the summed links into full-duration detector measurements.
signal_instrument = lisainstrument.Instrument(
    **{key: value for key, value in detector_config.items()
       if key not in ('T_obs', 'gpu_batch_size', 'minimum_envelope_samples', 'include_padding')}, gws=str(links_file), size=N + 2 * N_day,
    t0=-N_day / fs,
)
signal_instrument.disable_all_noises()
signal_instrument.disable_dopplers()
signal_data = pytdi.Data.from_instrument(signal_instrument.export_numpy_full())
signal_measurements = signal_data.measurements
noise_measurements = noise_data.measurements
del noise_data, noise_instrument

# Step 5. Combine the full observation, apply TDI, and save.
total_measurements = {
    name: signal_measurements[name] + noise_measurements[name]
    for name in signal_measurements
}
tdi_results = {}
for name, measurements in (('total', total_measurements),
                           ('signal', signal_measurements), ('noise', noise_measurements)):
    tdi_results[name] = tdi_gpu(
        measurements, signal_data.args, keep=slice(N_day, N_day + N),
    ) / signal_instrument.central_freq
detector_config['orbits'] = str(observation_file.with_suffix('.orbits.h5'))
with h5py.File(workspace / 'population_tdi.h5', 'x') as saved:
    saved.attrs.update(duration=T_obs, rate=fs)
    saved.attrs['detector_config'] = json.dumps(detector_config)
    for name, value in tdi_results.items():
        saved.create_dataset(name + '_xyz_td', data=value, chunks=(3, min(N_day, N)))

pack_observation(workspace, observation_file, overwrite=True, time_domain_only=True)
shutil.copy2(workspace / 'orbit.h5', observation_file.with_suffix('.orbits.h5'))
shutil.rmtree(workspace)
print(f'Saved {source_count:,} DWDs: {observation_file.resolve()}', flush=True)
\end{lstlisting}

\Needspace{8\baselineskip}
\subsection{Measured covariance and template F-statistics: script/dwd\_blind\_search.py}


% SOURCE script/dwd_blind_search.py 1 68
\begin{lstlisting}[language=Python,firstnumber=1,basicstyle=\ttfamily\scriptsize,numbers=left]
"""Measured covariance and analytical F-statistics for an existing DWD template bank."""
import numpy as np


def residual_covariance(data, duration, rate, *, neighboring_bins=256, chunk_size=16384):
    """Estimate the measured XYZ covariance without rejecting any frequency blocks."""
    from script.eric_preprocess import effective_noise_cov, periodogram

    data = np.asarray(data)
    frequencies = np.arange(data.shape[-1]) / duration
    covariance = np.zeros((len(frequencies), 3, 3), complex)
    for first in range(1, len(frequencies), chunk_size):
        last = min(first + chunk_size, len(frequencies))
        lower, upper = max(1, first - neighboring_bins), min(len(frequencies), last + neighboring_bins)
        psd = periodogram(data[:, lower:upper], duration, rate)
        if not np.any(psd):
            local = np.broadcast_to(np.eye(3) * 1.e-100, psd.shape).copy()
        else:
            local = effective_noise_cov(psd, neighboring_bins=neighboring_bins, downweight_peaks=True)
            scale = np.trace(local, axis1=-2, axis2=-1).real / 3
            floor = max(float(np.median(scale)) * 1.e-10, 1.e-100)
            local += floor * np.eye(3)
        retained = local[first - lower:last - lower]
        covariance[first:last] = retained
    covariance[0] = covariance[1]
    return covariance


def amplitude_f_statistic(data, bases, covariance, duration, rate):
    """Vectorized standard real four-amplitude projection, including XYZ cross terms."""
    bases = np.asarray(bases)
    inverse = np.linalg.inv(covariance)
    factor = 4 / (duration * rate**2)
    x = factor * np.einsum('smcf,fcd,df->sm', bases.conj(), inverse, data, optimize=True).real
    M = factor * np.einsum('smcf,fcd,sndf->smn', bases.conj(), inverse, bases, optimize=True).real
    M = .5 * (M + M.swapaxes(-1, -2))
    # A rank-deficient basis still defines a projection (especially near Nyquist).
    coefficients = np.einsum('smn,sn->sm', np.linalg.pinv(M, hermitian=True, rcond=1.e-12), x)
    two_f = np.maximum(np.einsum('sm,sm->s', x, coefficients), 0.)
    return dict(two_f=two_f, coefficients=coefficients, amplitude_matrix=M)


def template_f_statistic_search(data, covariance, builder, core_indices, duration, rate,
                                template_parameters, *, f_threshold=16., batch_size=128):
    """Evaluate F at every supplied template; maximize its four amplitudes analytically."""
    templates = np.asanyarray(template_parameters)
    if templates.dtype.names:
        doppler = np.column_stack([templates[name] for name in
            ('Frequency', 'FrequencyDerivative', 'EclipticLongitude', 'EclipticLatitude')])
    elif templates.ndim == 2 and templates.shape[1] == 11:
        doppler = templates[:, [2, 3, 6, 7]]
    else:
        templates = templates.reshape(-1, 4)
        doppler = templates
    scores = np.empty(len(templates))
    coefficients = np.empty((len(templates), 4))
    for first in range(0, len(templates), batch_size):
        points = doppler[first:first + batch_size]
        padded = np.vstack((points, np.repeat(points[-1:], batch_size - len(points), axis=0)))
        bases = builder.build(padded)[:len(points), :, :, core_indices]
        fit = amplitude_f_statistic(data, bases, covariance, duration, rate)
        scores[first:first + len(points)] = fit['two_f']
        coefficients[first:first + len(points)] = fit['coefficients']
    retained = scores >= 2 * f_threshold
    return dict(template_dwd=templates[retained], doppler=doppler[retained], two_f=scores[retained],
        coefficients=coefficients[retained], template_parameters=templates,
        template_two_f=scores, retained_template_indices=np.flatnonzero(retained),
        template_evaluations=len(templates))
\end{lstlisting}

\Needspace{8\baselineskip}
\subsection{DWD synthesis and native-expression GPU TDI: eric\_generate\_waveform.py}


% SOURCE eric_generate_waveform.py 1 266
\begin{lstlisting}[language=Python,firstnumber=1,basicstyle=\ttfamily\scriptsize,numbers=left,lineskip=-0.2pt]
"""GPU DWD link synthesis and native LISA time-delay interferometry."""
import numpy as np


def DWD_simulation_gpu(DWD_parameters, detector_config):
    """Return the summed six-link response of the supplied DWD parameter rows.

    DWD_parameters is a float array with shape (n_DWD, 8), with columns:
    A, f, df, phi0, iota, ecliptic longitude, ecliptic latitude, polarization.
    An (n_DWD, 11) SNR table is also accepted: ID, those eight parameters,
    SNR_absolute, SNR_relative. The ID and SNR columns are ignored.
    Alternatively, pass a one-dimensional structured table with fields:
    Amplitude, Frequency, FrequencyDerivative, InitialPhase, Inclination,
    EclipticLongitude, EclipticLatitude, Polarization.
    Additional fields such as ID and SNR are ignored. Every supplied row
    contributes once, including repeated parameter sets. No catalogue or ID
    lookup is performed; sources need not appear in an LDC file.

    detector_config requires T_obs [s], dt [s], and orbits (orbit-file path).
    The equivalent fs [Hz] and orbit_file keys remain supported. Arm lengths
    and spacecraft motion come from that orbit file. Optional settings are
    gpu_batch_size (default 10000) and minimum_envelope_samples (default 512).
    include_padding=True retains two days on each side for instrument/TDI use.

    By default, the result is a float64 NumPy array with shape (round(T_obs * fs), 6),
    sampled at t = arange(N) / fs in native LISA link order
    (12, 23, 31, 13, 32, 21). Values are dimensionless link fractional
    frequency shifts, before instrument and TDI processing. No noise is added
    and no output file is written. CUDA is required for a nonempty table.
    With include_padding=True, N_pad = round(2 * 86400 * fs) samples are retained
    at each end; the first sample is at t = -N_pad / fs.
    """
    T_obs = float(detector_config['T_obs'])
    fs = float(detector_config['fs']) if 'fs' in detector_config else 1 / float(detector_config['dt'])
    N = int(round(T_obs * fs))
    include_padding = bool(detector_config.get('include_padding', False))
    padding = int(round(2 * 86400 * fs)) if include_padding else 0
    batch_size = int(detector_config.get('gpu_batch_size', 10000))
    DWD_parameters = np.asarray(DWD_parameters)
    fields = ('Amplitude', 'Frequency', 'FrequencyDerivative', 'InitialPhase',
              'Inclination', 'EclipticLongitude', 'EclipticLatitude', 'Polarization')
    named = DWD_parameters.dtype.names is not None
    if len(DWD_parameters) == 0:
        return np.zeros((N + 2 * padding, 6), dtype=np.float64)

    from astropy import units as u
    from astropy.coordinates import BarycentricMeanEcliptic, CartesianRepresentation, SkyCoord
    from astropy.time import Time
    from script.dwd_population_gpu import PopulationLinksGPU
    # Rotate the supplied sky position and polarization basis into ICRS.
    rotation = np.empty((3, 3))
    frame = BarycentricMeanEcliptic(equinox=Time('J2000'))
    for column, axis in enumerate(np.eye(3)):
        rotation[:, column] = SkyCoord(
            CartesianRepresentation(*(axis * u.one)), frame=frame,
        ).icrs.cartesian.xyz.value
    engine = PopulationLinksGPU(
        (detector_config['orbit_file'] if 'orbit_file' in detector_config else detector_config['orbits']), T_obs, fs, batch_size=batch_size,
        minimum=int(detector_config.get('minimum_envelope_samples', 512)),
    )
    for first in range(0, len(DWD_parameters), 2 * batch_size):
        selected = DWD_parameters[first:first + 2 * batch_size]
        values = (np.column_stack([selected[field] for field in fields])
                  if named else np.asarray(selected, dtype=np.float64))
        if not named and values.shape[1] == 11:
            values = values[:, 1:9]
        longitude, latitude, polarization = values[:, 5], values[:, 6], values[:, 7]
        cos_lon, sin_lon = np.cos(longitude), np.sin(longitude)
        cos_lat, sin_lat = np.cos(latitude), np.sin(latitude)
        direction = np.column_stack((cos_lat * cos_lon, cos_lat * sin_lon, sin_lat)) @ rotation.T
        transported_u = np.column_stack((sin_lon, -cos_lon, np.zeros_like(longitude))) @ rotation.T
        ra = np.mod(np.arctan2(direction[:, 1], direction[:, 0]), 2 * np.pi)
        dec = np.arcsin(np.clip(direction[:, 2], -1., 1.))
        u_icrs = np.column_stack((np.sin(ra), -np.cos(ra), np.zeros_like(ra)))
        v_icrs = np.column_stack((-np.sin(dec) * np.cos(ra),
                                  -np.sin(dec) * np.sin(ra), np.cos(dec)))
        basis_rotation = np.arctan2(np.einsum('ij,ij->i', transported_u, v_icrs),
                                   np.einsum('ij,ij->i', transported_u, u_icrs))
        psi = np.mod(polarization + basis_rotation, np.pi)
        engine.add(np.column_stack((values[:, :5], ra, dec, psi)))
    response = engine.finish()
    return (response if include_padding else response[engine.padding:engine.padding + N]).copy()


# Lagrange interpolation adapted from PyTDI (BSD 3-Clause).
# Copyright (c) 2022, California Institute of Technology and
# Max Planck Institute for Gravitational Physics (Albert Einstein Institute).
#
# Redistribution and use in source and binary forms, with or without
# modification, are permitted provided that the following conditions are met:
# 1. Redistributions of source code must retain the above copyright notice,
#    this list of conditions and the following disclaimer.
# 2. Redistributions in binary form must reproduce the above copyright notice,
#    this list of conditions and the following disclaimer in the documentation
#    and/or other materials provided with the distribution.
# 3. Neither the name of the copyright holder nor the names of its contributors
#    may be used to endorse or promote products derived from this software
#    without specific prior written permission.
#
# THIS SOFTWARE IS PROVIDED BY THE COPYRIGHT HOLDERS AND CONTRIBUTORS "AS IS"
# AND ANY EXPRESS OR IMPLIED WARRANTIES, INCLUDING, BUT NOT LIMITED TO, THE
# IMPLIED WARRANTIES OF MERCHANTABILITY AND FITNESS FOR A PARTICULAR PURPOSE
# ARE DISCLAIMED. IN NO EVENT SHALL THE COPYRIGHT HOLDER OR CONTRIBUTORS BE
# LIABLE FOR ANY DIRECT, INDIRECT, INCIDENTAL, SPECIAL, EXEMPLARY, OR
# CONSEQUENTIAL DAMAGES (INCLUDING, BUT NOT LIMITED TO, PROCUREMENT OF
# SUBSTITUTE GOODS OR SERVICES; LOSS OF USE, DATA, OR PROFITS; OR BUSINESS
# INTERRUPTION) HOWEVER CAUSED AND ON ANY THEORY OF LIABILITY, WHETHER IN
# CONTRACT, STRICT LIABILITY, OR TORT (INCLUDING NEGLIGENCE OR OTHERWISE)
# ARISING IN ANY WAY OUT OF THE USE OF THIS SOFTWARE, EVEN IF ADVISED OF THE
# POSSIBILITY OF SUCH DAMAGE.

_CUDA_SHIFT = r'''
extern "C" __global__ void shift(
    const double *data, const double *shifts, const double *ratios,
    double *result, const long n, const int batches, const int halfp) {
    long i=(long)blockDim.x*blockIdx.x+threadIdx.x;
    if (i>=n) return;
    double q=floor(shifts[i]);
    double f=__dsub_rn(shifts[i],q);
    double taps[32];
    double factor=__dmul_rn(f,__dsub_rn(1.,f));
    for (int j=1;j<halfp;++j) {
        factor=__dmul_rn(factor,ratios[j]);
        taps[halfp-1-j]=__ddiv_rn(factor,__dadd_rn((double)j,f));
        taps[halfp+j]=__ddiv_rn(factor,__dsub_rn((double)(j+1),f));
    }
    taps[halfp-1]=__dsub_rn(1.,f);
    taps[halfp]=f;
    for (int j=2;j<halfp;++j) {
        double a=__ddiv_rn(f,(double)j);
        double b=__dsub_rn(1.,__dmul_rn(a,a));
        for (int k=0;k<2*halfp;++k) taps[k]=__dmul_rn(taps[k],b);
    }
    if (halfp>1) {
        double b=__dmul_rn(__dadd_rn(1.,f),__dsub_rn(1.,__ddiv_rn(f,(double)halfp)));
        for (int k=0;k<2*halfp;++k) taps[k]=__dmul_rn(taps[k],b);
    }
    long center=i+(long)q;
    if (center<-(halfp+1)) center=-(halfp+1);
    if (center>n+halfp-1) center=n+halfp-1;
    for (int batch=0;batch<batches;++batch) {
        double value=0.;
        for (int k=0;k<2*halfp;++k) {
            long index=center-(halfp-1)+k;
            double sample=(index>=0 && index<n)?data[(long)batch*n+index]:0.;
            value=__dadd_rn(value,__dmul_rn(taps[k],sample));
        }
        result[(long)batch*n+i]=value;
    }
}
'''


_TDI_ORBITS = {}


def tdi_gpu(mosa_links, detector_config, combination='XYZ2', *, t0=0., keep=slice(None)):
    """Return native-LISA XYZ time streams from six MOSA link measurements.

    mosa_links is an array of shape (..., N, 6), in native link order
    (12, 23, 31, 13, 32, 21). These are science-link measurements, with
    any instrument filtering already applied. The output is float64 with
    shape (..., 3, N), in X/Y/Z order and in the same units as the input.
    No FFT, instrument simulation, or file output is performed here.

    A dictionary of native PyTDI measurements is also accepted. Include
    sci_ij, tmi_ij and ref_ij when using full instrument noise: keeping these
    auxiliary beatnotes preserves the standard noise combinations. Missing
    beatnotes are zero, as in native Data.from_orbits. Keys '12', ..., '21'
    are shorthand for the six sci_ij measurements. Leading axes batch
    independent streams; the final axis of each dictionary value is time.

    detector_config requires fs [Hz] or dt [s], and orbits (or orbit_file).
    Moving unequal-arm light-travel times and derivatives come from that
    file, using orbit_dataset='tcb/ltt' by default. t0 gives the first
    input sample time. Native Data.args (fs, delays, delay_derivatives) can
    be passed instead to reuse already evaluated moving-orbit geometry.
    combination='XYZ2' selects second-generation Michelson by default;
    'XYZ1' selects first generation. keep is a slice removing input padding.
    Native orders 5/31 and FP64 arithmetic are preserved; CUDA is required.
    """
    import ctypes
    import copy
    from collections.abc import Mapping
    from pathlib import Path
    import sys
    from types import FunctionType, MethodType, SimpleNamespace
    import pytdi
    import pytdi.michelson

    channel_names = tuple(axis + combination[-1] for axis in 'XYZ')
    if isinstance(mosa_links, Mapping):
        provided = {('sci_' + str(name) if str(name) in pytdi.Data.MOSAS else str(name)):
                    np.asarray(value, dtype=np.float64) for name, value in mosa_links.items()}
    else:
        values = np.asarray(mosa_links, dtype=np.float64)
        provided = {'sci_' + mosa: values[..., column]
                    for column, mosa in enumerate(pytdi.Data.MOSAS)}
    shapes = [value.shape for value in provided.values() if value.ndim]
    shape = np.broadcast_shapes(*shapes)
    size = shape[-1]
    fs = float(detector_config['fs']) if 'fs' in detector_config else 1 / float(detector_config['dt'])
    if 'delays' in detector_config:
        native_args = detector_config
    else:
        orbit = Path(detector_config.get('orbit_file', detector_config.get('orbits', ''))).expanduser().resolve()
        stat = orbit.stat()
        dataset = detector_config.get('orbit_dataset', 'tcb/ltt')
        key = (str(orbit), stat.st_size, stat.st_mtime_ns, size, fs, float(t0), dataset)
        if key not in _TDI_ORBITS:
            data = pytdi.Data.from_orbits(str(orbit), fs=fs, t0=float(t0), dataset=dataset,
                                         sci_12=np.zeros(size, dtype=np.float64))
            _TDI_ORBITS.clear()
            _TDI_ORBITS[key] = data.args
        native_args = _TDI_ORBITS[key]

    cuda_root = Path(sys.prefix) / 'lib' / f'python{sys.version_info.major}.{sys.version_info.minor}' / 'site-packages/nvidia'
    handles = [ctypes.CDLL(str(cuda_root / package / 'lib' / library), mode=ctypes.RTLD_GLOBAL)
               for package, library in (('cuda_runtime', 'libcudart.so.12'), ('cuda_nvrtc', 'libnvrtc.so.12'))]
    import cupy as cp
    kernel = cp.RawKernel(_CUDA_SHIFT, 'shift', options=('--fmad=false',))
    ratios = {halfp: cp.asarray(np.array([(-1)*(1-j/halfp)/(1+j/halfp)
                                          for j in range(halfp)], dtype=np.float64))
              for halfp in (3, 16)}

    def timeshift(values, shifts, order=31):
        if np.isscalar(values):
            return values
        if np.isscalar(shifts) and shifts == 0:
            return values
        values = cp.ascontiguousarray(values, dtype=cp.float64)
        shifts = cp.asarray(shifts, dtype=cp.float64)
        if values.size == 1:
            return values.reshape(())
        if shifts.size != size:
            shifts = cp.broadcast_to(shifts, (size,)).copy()
        halfp = (order + 1) // 2
        result = cp.empty_like(values)
        batches = values.size // size
        kernel(((size + 127)//128,), (128,),
               (values, shifts, ratios[halfp], result, np.int64(size), np.int32(batches), np.int32(halfp)))
        return result

    gpu_dsp = SimpleNamespace(timeshift=timeshift, diff=pytdi.dsp.diff)
    methods = {}
    for name in ('build_shifts', 'build'):
        original = getattr(pytdi.TDICombination, name)
        method = FunctionType(original.__code__, {**original.__globals__, 'dsp': gpu_dsp},
                              original.__name__, original.__defaults__, original.__closure__)
        method.__kwdefaults__ = original.__kwdefaults__
        methods[name] = method
    args = {'delays': {name: cp.asarray(value, dtype=cp.float64) for name, value in native_args['delays'].items()},
            'delay_derivatives': {name: cp.asarray(value, dtype=cp.float64) for name, value in native_args['delay_derivatives'].items()},
            'fs': fs}
    names = tuple(dict.fromkeys(name for channel in channel_names
                               for name in getattr(pytdi.michelson, channel).components))
    measurements = {name: cp.asarray(np.ascontiguousarray(
        np.broadcast_to(provided.get(name, 0.), shape).reshape(-1, size))) for name in names}
    channels = []
    for name in channel_names:
        native_combination = copy.copy(getattr(pytdi.michelson, name))
        native_combination.build_shifts = MethodType(methods['build_shifts'], native_combination)
        operator = methods['build'](native_combination, **args)
        channels.append(cp.asnumpy(operator(measurements)[..., keep]))
    xyz = np.stack(channels, axis=-2)
    return xyz.reshape(*shape[:-1], 3, xyz.shape[-1])
\end{lstlisting}

\Needspace{8\baselineskip}
\subsection{Variable-count inference: eric\_rjmcmc.py}


% SOURCE eric_rjmcmc.py 1 316
\begin{lstlisting}[language=Python,firstnumber=1,basicstyle=\ttfamily\scriptsize,lineskip=-0.2pt,numbers=left]
"""DWD priors, amplitude-marginal likelihood, and GPU-only Eryn sampling."""

from functools import lru_cache
from importlib.metadata import version
from pathlib import Path
import hashlib
import time

import jax
import jax.numpy as jnp
import jax.scipy.linalg as jsl
import numpy as np
from scipy.linalg import cho_factor, cho_solve, solve_triangular
from scipy.special import comb, gammaln
from tqdm.auto import tqdm
from eryn.backends import HDFBackend
from eryn.ensemble import EnsembleSampler
from eryn.moves import DistributionGenerateRJ, GaussianMove
from eryn.prior import ProbDistContainer
from eryn.state import State

from script.eric_generate_waveform import GPUResponse
from script.eric_preprocess import inner_product
from script.tool import single_cpu


# 1. Priors and sampled coordinates
def encode_doppler(para_doppler, f_center, T_obs, *, chirp_coordinate='log10'):
    para_doppler = np.asarray(para_doppler)
    if para_doppler.dtype.names:
        para_doppler = np.column_stack([para_doppler[name] for name in
            ('Frequency', 'FrequencyDerivative', 'EclipticLongitude', 'EclipticLatitude')])
    elif para_doppler.ndim == 2 and para_doppler.shape[1] == 11:
        para_doppler = para_doppler[:, [2, 3, 6, 7]]
    else:
        para_doppler = para_doppler.reshape(-1, 4)
    chirp = (para_doppler[:, 1]*T_obs**2 if chirp_coordinate == 'linear'
             else np.log10(para_doppler[:, 1]))
    return np.column_stack((np.log10(para_doppler[:, 0]/f_center)*f_center*T_obs*np.log(10.),
                            chirp, para_doppler[:, 2] % (2*np.pi),
                            np.sin(para_doppler[:, 3])))


def decode_doppler(para_doppler, f_center, T_obs, *, chirp_coordinate='log10'):
    para_doppler = np.asarray(para_doppler).reshape(-1, 4)
    chirp = (para_doppler[:, 1]/T_obs**2 if chirp_coordinate == 'linear'
             else 10.**para_doppler[:, 1])
    return np.column_stack((f_center*10.**(para_doppler[:, 0]/(f_center*T_obs*np.log(10.))),
                            chirp, para_doppler[:, 2] % (2*np.pi),
                            np.arcsin(np.clip(para_doppler[:, 3], -1., 1.))))


class CataloguePrior:
    def __init__(self, bounds, maximum=49):
        self.bounds, self.count_maximum = np.asarray(bounds), maximum
        self.V = np.prod(self.bounds[:, 1] - self.bounds[:, 0])
        self.DWD_para_order = [0, 1, 2, 3]
        self.key_order = self.DWD_para_order

    def logpdf(self, coords, inds, **kwargs):
        para_doppler, active = coords['dwd'], inds['dwd']
        N = active.sum(axis=-1)
        M = self.count_maximum
        V = self.V
        value = -N * np.log(V) - np.log(comb(M, N))
        outside = np.any((para_doppler < self.bounds[:, 0]) | (para_doppler > self.bounds[:, 1])
                         | ~np.isfinite(para_doppler), axis=-1) & active
        return np.where(outside.any(axis=-1), -np.inf, value)


class BirthDistribution:
    """90% initial-candidate boxes; 10% full-prior support."""
    def __init__(self, centers, bounds):
        self.centers, self.bounds = centers, bounds
        self.lower_limit = np.maximum(centers-[.2, 0., .2, .2], bounds[:, 0])
        self.upper_limit = np.minimum(centers+[.2, 0., .2, .2], bounds[:, 1])
        self.lower_limit[:, 1], self.upper_limit[:, 1] = bounds[1]
        self.lower_limit[:, 2], self.upper_limit[:, 2] = centers[:, 2]-.2, centers[:, 2]+.2

    def logpdf(self, para_doppler):
        para_doppler = np.asarray(para_doppler).reshape(-1, 4)
        broad = 1./np.prod(self.bounds[:, 1]-self.bounds[:, 0])
        if not len(self.centers):
            return np.full(len(para_doppler), np.log(broad))
        delta = (para_doppler[:, None, 2]-self.centers[:, 2]+np.pi) % (2*np.pi)-np.pi
        inside = np.all((para_doppler[:, None, [0, 1, 3]] >= self.lower_limit[:, [0, 1, 3]])
                        & (para_doppler[:, None, [0, 1, 3]] <= self.upper_limit[:, [0, 1, 3]]), axis=-1)
        inside &= np.abs(delta) <= .5*(self.upper_limit[:, 2]-self.lower_limit[:, 2])
        density = .1*broad+.9*(inside/np.prod(self.upper_limit-self.lower_limit, axis=1)).mean(axis=-1)
        return np.log(density)

    def rvs(self, size=1):
        shape = (size,) if isinstance(size, int) else tuple(size)
        para_doppler = np.random.uniform(self.bounds[:, 0], self.bounds[:, 1], size=shape+(4,))
        if len(self.centers):
            flat = para_doppler.reshape(-1, 4)
            selected = np.flatnonzero(np.random.random(len(flat)) >= .1)
            candidate = np.random.randint(len(self.centers), size=len(selected))
            flat[selected] = np.random.uniform(self.lower_limit[candidate], self.upper_limit[candidate])
        para_doppler[..., 2] %= 2*np.pi
        return para_doppler


# 2. Likelihood
def marginal_likelihood(data, bases, noise_covariance, log_determinant, T_obs, rate,
                        amplitude_prior_sigma=10.0):
    """Full 4N Gaussian integral using the noise covariance, not its inverse."""
    data_norm = inner_product(data, data, noise_covariance, T_obs, rate)
    if len(bases):
        h = bases.reshape(len(bases)*4, *data.shape)
        x = np.asarray([inner_product(a, data, noise_covariance, T_obs, rate) for a in h])
        M = np.asarray([[inner_product(a, b, noise_covariance, T_obs, rate)
                         for b in h] for a in h])
        M = 0.5 * (M + M.T)
        precision = M + np.eye(len(M)) / amplitude_prior_sigma**2
        chol, lower = cho_factor(precision, lower=True, check_finite=False)
        mean = cho_solve((chol, lower), x, check_finite=False)
        covariance = cho_solve((chol, lower), np.eye(len(M)), check_finite=False)
        penalty = np.log(np.diag(chol)).sum() + len(M)*np.log(amplitude_prior_sigma)
        improvement = 0.5 * x @ mean - penalty
        coefficients = mean.reshape(-1, 4)
        model = np.einsum('am,amck->ck', coefficients, bases, optimize=True)
    else:
        M = covariance = np.empty((0, 0))
        x = np.empty(0)
        coefficients = np.empty((0, 4))
        improvement = penalty = 0.0
        model = np.zeros_like(data)
    residual = data - model
    residual_norm = inner_product(residual, residual, noise_covariance, T_obs, rate)
    return dict(log_marginal=float(-0.5*data_norm-log_determinant+improvement),
                log_likelihood=float(-0.5*residual_norm-log_determinant),
                coefficients=coefficients, amplitude_covariance=covariance,
                projection=x, data_norm=float(data_norm), log_determinant=float(log_determinant),
                amplitude_matrix=M, log_occam_penalty=float(penalty),
                reconstructed=model, residual=residual, residual_norm=float(residual_norm))


class Likelihood:
    def __init__(self, engine, f_center, T_obs, sigma=10., *, maximum_sources=49, chirp_coordinate='log10'):
        self.f_center, self.T_obs, self.sigma = f_center, T_obs, sigma
        self.chirp_coordinate = chirp_coordinate
        self.engine, self.count_maximum = engine, maximum_sources
        self.white_data = engine.white_data
        self.data_norm = float(self.white_data @ self.white_data)
        self.waveform_evaluations, self.calls = 0, 0
        self.source = lru_cache(maxsize=256)(self._source)
        self.cached_details = lru_cache(maxsize=64)(self._details)
        self.integral = jax.jit(self._integral)
        self.zero = jax.device_put(np.zeros((4, self.white_data.size)), engine.device)

    def _source(self, key):
        para_doppler = np.frombuffer(key, dtype=np.float64)
        point = decode_doppler(para_doppler[None], self.f_center, self.T_obs,
                               chirp_coordinate=self.chirp_coordinate)[0]
        self.waveform_evaluations += 1
        return self.engine.white_single(jax.device_put(point, self.engine.device),
                                        *self.engine.arguments)

    def _integral(self, white):
        gram = white @ white.T
        lower = jnp.linalg.cholesky(gram+jnp.eye(len(white))/self.sigma**2)
        x = white @ self.white_data
        A = jsl.cho_solve((lower, True), x)
        logm = (-.5*self.data_norm+.5*x@A
                -jnp.log(jnp.diag(lower)).sum()-len(white)*jnp.log(self.sigma))
        return logm, A, lower

    def _details(self, key):
        para_doppler = np.frombuffer(key, dtype=np.float64).reshape(-1, 4)
        if not len(para_doppler):
            return -.5*self.data_norm, np.empty(0), np.empty((0, 0))
        white = [self.source(point.tobytes()) for point in para_doppler]
        white += [self.zero]*(self.count_maximum-len(para_doppler))
        logm, A, lower = self.integral(jnp.concatenate(white, axis=0))
        return float(logm), np.asarray(A)[:4*len(para_doppler)], np.asarray(lower)[:4*len(para_doppler), :4*len(para_doppler)]

    def details(self, para_doppler):
        return self.cached_details(np.ascontiguousarray(para_doppler, dtype=np.float64).tobytes())

    def __call__(self, para_doppler):
        self.calls += 1
        return self.details(para_doppler)[0]


# 3. Eryn samplers
class BirthDeathMove(DistributionGenerateRJ):
    def get_proposal(self, coords, inds, *args, **kwargs):
        para_doppler, new_inds, factors = super().get_proposal(coords, inds, *args, **kwargs)
        n, new_n = inds['dwd'].sum(axis=-1), new_inds['dwd'].sum(axis=-1)
        maximum = inds['dwd'].shape[-1]
        # Slot-selection ratio paired with CataloguePrior's uniform count prior.
        factors += (gammaln(n+1)+gammaln(maximum-n+1)
                    -gammaln(new_n+1)-gammaln(maximum-new_n+1))
        return para_doppler, new_inds, factors


class DWDMove(GaussianMove):
    def get_proposal(self, coords, random, branches_inds=None, **kwargs):
        active = branches_inds['dwd']
        slot = np.argmax(np.where(active, random.random(active.shape), -1.), axis=-1)
        selected = np.zeros_like(active)
        np.put_along_axis(selected, slot[..., None], active.any(axis=-1)[..., None], axis=-1)
        para_doppler, factors = super().get_proposal(coords, random, branches_inds={'dwd': selected})
        beta = self.temperature_control.betas[:, None, None, None]
        para_doppler['dwd'] = coords['dwd']+(para_doppler['dwd']-coords['dwd'])/np.sqrt(beta)
        para_doppler['dwd'][..., 2] %= 2*np.pi
        return para_doppler, factors


def run_eryn_rjmcmc(data, covariance, bank_doppler,
        frequency_range, derivative_bounds, f_center, T_obs, rate,
        checkpoint_directory, *, response_times, response_geometry, physical_frequency,
        chain_count=1, inverse_temperatures=(1., .8, .6, .4),
        iterations=6000, burn_in=1000, maximum_sources=49, amplitude_prior_sigma=10.,
        chirp_coordinate='log10'):
    """Fit one residual with fixed covariance on GPU; the caller owns outer rounds."""
    with single_cpu():
        started = time.perf_counter()
        maximum, sigma = maximum_sources, amplitude_prior_sigma
        betas = np.asarray(inverse_temperatures)
        checkpoint_file = Path(checkpoint_directory)/'eryn_chain_00.h5'
        covariance = np.array(covariance, dtype=np.complex128, copy=True)
        covariance_sha = hashlib.sha256(covariance.tobytes()).hexdigest()
        scale = f_center*T_obs*np.log(10.)
        chirp_bounds = np.asarray(derivative_bounds, dtype=float)
        chirp_bounds = chirp_bounds*T_obs**2 if chirp_coordinate == 'linear' else np.log10(chirp_bounds)
        bounds = np.array([np.log10(np.asarray(frequency_range)/f_center)*scale,
                           chirp_bounds, [0., 2*np.pi], [-1., 1.]])
        centers = encode_doppler(bank_doppler, f_center, T_obs, chirp_coordinate=chirp_coordinate)

        engine = GPUResponse(covariance, data, response_times, response_geometry,
                             physical_frequency, T_obs, rate)
        likelihood = Likelihood(engine, f_center, T_obs, sigma, maximum_sources=maximum,
                                chirp_coordinate=chirp_coordinate)
        checkpoint_file.parent.mkdir(parents=True, exist_ok=True)
        prior = CataloguePrior(bounds, maximum)
        within = DWDMove({'dwd': np.diag(np.array([.025, .3, .015, .015])**2)})
        births = ProbDistContainer({(0, 1, 2, 3): BirthDistribution(centers, bounds)})
        jump = BirthDeathMove({'dwd': births}, nleaves_min={'dwd': 0},
                             nleaves_max={'dwd': maximum}, prevent_swaps=True)
        backend = HDFBackend(str(checkpoint_file))
        sampler = EnsembleSampler(
            1, {'dwd': 4}, likelihood, {'dwd': prior, 'all_models_together': prior},
            branch_names=['dwd'], nleaves_min={'dwd': 0}, nleaves_max={'dwd': maximum},
            moves=within, rj_moves=jump, pool=None, backend=backend,
            tempering_kwargs=dict(betas=betas, adaptive=False, permute=False),
            fill_zero_leaves_val=likelihood(np.empty((0, 4))),
        )
        para_doppler = np.zeros((len(betas), 1, maximum, 4))
        active = np.zeros(para_doppler.shape[:-1], dtype=bool)
        para_doppler[:, 0, :len(centers)], active[:, 0, :len(centers)] = centers, True
        state = State({'dwd': para_doppler}, inds={'dwd': active})
        initial_counts = active[:, 0].sum(axis=-1)
        initial_logm = np.full(len(betas), likelihood(centers))
        initial_prior = prior.logpdf(state.branches_coords, state.branches_inds)[:, 0]
        setup_seconds = time.perf_counter()-started

        para_doppler_draws, a_draws, count_draws, logm_draws = [], [], [], []
        amplitude_rng = np.random.default_rng()
        progress = tqdm(sampler.sample(state, iterations=iterations), total=iterations,
                        desc='Eryn RJMCMC', unit='cycle', dynamic_ncols=True)
        for cycle, state in enumerate(progress, 1):
            progress.set_postfix(N=state.branches_inds['dwd'][:, 0].sum(axis=-1).tolist(),
                                 refresh=False)
            if cycle <= burn_in:
                continue
            para_doppler = state.branches_coords['dwd'][0, 0, state.branches_inds['dwd'][0, 0]]
            logm, A, lower = likelihood.details(para_doppler)
            draw = A+(solve_triangular(lower.T, amplitude_rng.normal(size=len(A)),
                                       lower=False) if len(A) else A)
            order = np.argsort(para_doppler[:, 0])
            padded_para_doppler, padded_a = np.full((maximum, 4), np.nan), np.full((maximum, 4), np.nan)
            padded_para_doppler[:len(para_doppler)], padded_a[:len(para_doppler)] = para_doppler[order], draw.reshape(-1, 4)[order]
            para_doppler_draws.append(padded_para_doppler)
            a_draws.append(padded_a)
            count_draws.append(len(para_doppler))
            logm_draws.append(logm)

        probability = np.bincount(count_draws, minlength=maximum+1)/len(count_draws)
        modal = int(np.argmax(probability))
        selected = np.flatnonzero(np.asarray(count_draws) == modal)
        best = selected[np.argmax(np.asarray(logm_draws)[selected])]
        doppler = decode_doppler(para_doppler_draws[best][:modal], f_center, T_obs,
                                 chirp_coordinate=chirp_coordinate)
        bases = engine.build(doppler)
        fit = marginal_likelihood(data, bases, covariance, 0., T_obs, rate, sigma)
        counts = np.vstack((initial_counts, backend.get_nleaves()['dwd'][:, :, 0]))
        logm = np.vstack((initial_logm, backend.get_log_like()[:, :, 0]))
        logp = np.vstack((initial_prior, backend.get_log_prior()[:, :, 0]))
        chain = dict(sampled_q=np.array(para_doppler_draws), sampled_amplitudes=np.array(a_draws),
            source_count=np.array(count_draws), log_marginal=np.array(logm_draws),
            trace_counts=counts[1:], trace_logm=logm[1:], initial_replica_counts=initial_counts,
            cycle_state=dict(cycle=np.arange(iterations+1),
                is_burn_in=np.arange(iterations+1) <= burn_in, source_count=counts,
                log_marginal=logm, log_prior=logp, log_posterior=logm*betas+logp))
        result = dict(doppler=doppler, bases=bases, fit_result=fit, modal_count=modal,
            count_probability=probability, count_support=np.arange(maximum+1), chains={'0': chain},
            candidate_centers=centers, initial_replica_counts=initial_counts,
            config=dict(bounds=bounds, maximum_sources=maximum, center=f_center, duration=T_obs,
                rate=rate, iterations=iterations, burn_in=burn_in, thin=1, chain_count=1,
                inverse_temperatures=betas, amplitude_prior_sigma=sigma,
                chirp_coordinate=chirp_coordinate,
                initialization='all_candidates', initial_selection=None,
                use_gpu=True, backend='gpu',
                independent_chain_convergence_assessed=False),
            engine='Eryn '+version('eryn'), checkpoint_file=str(checkpoint_file),
            fixed_covariance_sha256=covariance_sha, backend_setup_seconds=setup_seconds,
            runtime_seconds=time.perf_counter()-started,
            waveform_evaluations=likelihood.waveform_evaluations, likelihood_calls=likelihood.calls,
            count_prior='Uniform N=0..maximum_sources; exchangeable active slots.',
            amplitude_prior=f'Normalized zero-mean Gaussian sigma={sigma} on each A_mu/1e-22.',
            proposal='Eryn Gaussian MH and distribution RJ; single-source updates; 90% candidate-box births.',
            estimate='Best cold marginal likelihood at modal N; conditional-mean amplitudes.')
        result.update(gpu_device=str(engine.device), gpu_device_kind=engine.device.device_kind)
        return result
\end{lstlisting}

\Needspace{8\baselineskip}
\section{Review reproduction}
\label{app:reproduce}

Runtime validation guards and explicit exceptions have been removed. The
waveform, TDI and statistical calculations retain their valid-input operations.
The TeX embeds source snapshots and is compilable without accessing the server.

\begin{lstlisting}[language=bash,numbers=left,firstnumber=1]
python -m unittest discover -s tests -p test_subtraction_all.py
python -m unittest discover -s tests -p test_blind_search.py
python -m unittest discover -s tests -p test_tdi_subtraction.py
# In the server CUDA scientific environment, using the saved native fixture:
python tests/smoke_subtraction_range.py
\end{lstlisting}

The source-listing checker compares every marked TeX listing with its named
Python file and verifies complete primary-file and SNR-producer coverage, current source hashes,
unique section labels and resolved cross-references. The PDF is separately
compiled and rendered to check title, contents, code, tables, mathematical
sections and page transitions.

\end{document}

